\exercisesection{正则环}

\begin{enumerate}
\def\labelenumi{\arabic{enumi})}
\item
  设 A 为 Noether 局部环。为使 A 正则，必须且只需 \(\mathrm{di}_{\mathrm{A}}(\kappa_{\mathrm{A}})<+\infty\)。
\item
  设 A 为正则环，p 为 A 的一个素理想。证明 A 关于 p-进拓扑的完备化是整环。
\item
  设 A 为环，\(\mathbf{T} = (T_i)_{i\in I}\) 为一个有限的未定元族。证明下列条件等价：
\end{enumerate}

\begin{enumerate}
\def\labelenumi{(\roman{enumi})}
\item
  环 A 是正则的；
\item
  环 A{[}T{]} 是正则的；
\item
  环 A{[}{[}T{]}{]} 是正则的。
\end{enumerate}

\begin{enumerate}
\def\labelenumi{\arabic{enumi})}
\setcounter{enumi}{3}
\item
  设 A 为维数 \(\leqslant 2\) 的 Noether 环。证明：为使 A 正则，必须且只需每个有限生成 A-模的对偶都是投射的（即设 \(\mathfrak{p}\) 取遍 \(\Lambda\) 的素理想；若第二个条件满足，先证明 \(\mathrm{A}_{\mathfrak{p}}\) 在 \(\operatorname{prof}(A_{\mathfrak{p}}) \leqslant 1\) 时是正则的，再证明 \(\mathrm{dp}_{\mathrm{A}_{\mathfrak{p}}}(\kappa(\mathfrak{p})) = 2\) 在 \(\operatorname{prof}(A_{\mathfrak{p}}) = 2\) 时成立）。
\item
  设 A 为同调维数有限的 Noether 局部环；试由 A, X, 第~208 页习题 16 推出 A 正则这一事实的另一证明（注意有 dim \(_{\kappa_A}\) \(\mathfrak{m}_A/\mathfrak{m}_A^2 \leqslant \mathrm{dh}(A) \leqslant \mathrm{prof}(A)\)）。
\item
  设 A 为正则局部环，\((x_{1},\ldots,x_{n})\) 与 \((y_{1},\ldots,y_{n})\) 为 \(m_{A}\) 中元素的极大割序列，分别生成理想 x 与 y。假设 \(y \subset x\)，于是存在 A 的元素 \(a_{ij}\)，使得 \(y_{i} = \sum_{j} a_{ij} x_{j}\) 对 \(1 \leqslant i \leqslant n\) 成立。证明 \(\det(a_{ij})\) 在 A/y 中的类不依赖于所选取的矩阵 \((a_{ij})\)，且其零化子是 x 的像；特别地，有 \(\det(a_{ij}) \neq 0\)（证明典范同态 \(\operatorname{Ext}_{\mathrm{A}}^{n}(\mathrm{A}/x,\mathrm{A}) \to \operatorname{Ext}_{\mathrm{A}}^{n}(\Lambda/\mathfrak{y},\Lambda)\) 是单射，再利用 Koszul 复形证明它与由比率为 \(\det(a_{ij})\) 的乘法映射诱导的同态 A/x → A/y 等同）。
\end{enumerate}

设 A 为维数 3 的正则局部环，f 为 \(m_{A}\) 的一个非零元素。证明下列条件等价：

\begin{enumerate}
\def\labelenumi{(\roman{enumi})}
\item
  环 A/(f) 不是因子分解整环；
\item
  存在整数 \(n \geqslant 2\) 以及一个矩阵 \(X \in \mathbf{M}_n(\mathrm{A})\)，其元素属于 \(\mathfrak{m}_{\mathrm{A}}\)，使得 \(f = \det(\mathrm{X})\)。
\end{enumerate}

（在条件 (ii) 下，证明 A/(f) 中由 X 的子式 \(X^{11}, \ldots, X^{1n}\)（其中 \(X^{1p}\) 由删去第一行与指标为 p 的列而得）的类生成的理想高度为 1 但不是主理想。在条件 (i) 下，设 p 为 A 的一个高度 2 的素理想，其在 \(\Lambda/(f)\) 中的像不是主理想；由 §3 习题 11 推出 p 由矩阵 \(Y \in \mathbf{M}_{n,n-1}(\mathbf{A})\)（元素属于 \(m_{\Lambda}\)）的 n-1 阶子式生成。把 f 写成矩阵 \(\widetilde{Y}\)（由对 Y 添加一列 \(Y_{n}\) 而得）的行列式；若 \(Y_{n}\) 中有一个元素可逆，则化为 \(Y_{n}\) 属于 \(\Lambda^{n}\) 的典范基的情形。

¶ 8) 设 A 为正则局部环，\(\mathbf{x} = (x_1, \ldots, x_d)\) 为 A 的一个坐标系，\(\rho : A \to B\) 为使 B 成为有限生成 A-模的单射环同态。证明下列条件等价：

\begin{enumerate}
\def\labelenumi{(\roman{enumi})}
\item
  对每个整数 \(p\)，元素 \(\rho(x_1 \ldots x_d)^p\) 不属于由 \(\rho(x_1^{p+1}), \ldots, \rho(x_d^{p+1})\) 生成的 B 的理想；
\item
  A-模 \(\rho(A)\) 是 B 的直和因子。
\end{enumerate}

当 A 是 Q-代数时这些条件成立（§ 2 习题 3）。

为证 (ii)⇒(i)，仿照该处的证明。对逆蕴含，化为 A 为完备的情形，并对 \(p \geqslant 0\) 令 \(A_{p} = A/(x_{1}^{p+1}, \ldots, x_{d}^{p+1})\)，\(B_{p} = B/(x_{1}^{p+1}, \ldots, x_{d}^{p+1})\)。注意 \((x_{1} \ldots x_{d})^{p}\) 的类生成 \(A_{p}\) 的底座，由此推出同态 \(\rho_{p}: A_{p} \to B_{p}\)（由 \(\rho\) 诱导）是单射。证明 \(\rho_{p}\) 容许一个 \(A_{p}\)-线性收缩。最后利用 E, III, 第~60 页，例 II。

\begin{enumerate}
\def\labelenumi{\arabic{enumi})}
\setcounter{enumi}{8}
\tightlist
\item
  设 \(\rho: A \to B\) 为 Noether 局部环之间的一个局部同态，\(N\) 为一个有限生成的 Macaulay B-模。假设环 \(A\) 正则，且我们有 \(\dim_{B}(N) = \dim(A) + \dim_{B/\mathfrak{m}_{A}B}(N/\mathfrak{m}_{A}N)\)。证明 A-模 \(N\) 是平坦的（对 \(\dim(A)\) 作归纳论证；若 \(x \in \mathfrak{m}_{A} - \mathfrak{m}_{A}^{2}\)，将证明 \(x\) 是 \(N\)-正则的，且同态 A/xA → B/xB 与 (B/xB)-模 N/xN 满足与 ρ 和 N 相同的假设）。
\end{enumerate}

¶ 10) 设 A 为 Noether 环，M 为有限生成 A-模。对每个整数 \(n \geqslant 0\)，用 \(X_{n}\) 表示元素 \(\mathfrak{p}\)（属于 \(\operatorname{Spec}(A)\)，满足 \(\dim_{A_{\mathfrak{p}}}(M_{\mathfrak{p}}) - \operatorname{prof}_{A_{\mathfrak{p}}}(M_{\mathfrak{p}}) \geqslant n\)）所成的集合。设 \(k\) 为一个整数 \(\geqslant 0\)。

\begin{enumerate}
\def\labelenumi{\alph{enumi})}
\item
  为使 M 满足性质 \(S_{k}\)（§ 1，习题 8），必须且只需 codim( \(X_{n}\)，Supp(M)) \(\geqslant n + k\) 对每个 \(n \geqslant 0\) 成立。
\item
  假设 \(X_{n}\) 在 \(\operatorname{Spec}(A)\) 中对每个 \(n \geqslant 0\) 都是闭的。设 p 为 \(\operatorname{Supp}(M)\) 的一个元素。为使 \(A_{p}\)-模 \(M_{p}\) 满足性质 \(S_{k}\)（§ 1，习题 8），必须且只需 \(\operatorname{codim}(X_{n}^{\prime}, \operatorname{Supp}(M)) \geqslant n + k\) 对每个 \(n \geqslant 0\) 以及每个分支 \(X_{n}^{\prime}\)（属于 \(X_{n}\) 且包含 p）成立。
\item
  假设环 A 可表现。由 b) 推出：A 中使 \(M_{p}\) 满足性质 \(S_{k}\) 的素理想 p 所成的集合在 Spec(A) 中是开的。
\end{enumerate}

\begin{enumerate}
\def\labelenumi{\arabic{enumi})}
\setcounter{enumi}{10}
\tightlist
\item
  设 A 为环 \(\mathbf{Z}[\mathbf{X}]\)，其中 \(\mathbf{X} = (X_i)_{i\in I}\) 为一个有限的未定元族。设 J 为 A 的一个由单项式生成的理想。
\end{enumerate}

\begin{enumerate}
\def\labelenumi{\alph{enumi})}
\item
  设 \(i \in I\)；设生成 J 的单项式为 \(\mathrm{P}_{1}, \ldots, \mathrm{P}_{r}\) 与 \(\mathrm{Q}_{1}, \ldots, \mathrm{Q}_{s}\)，其中 \(\mathrm{P}_{j}\) 被 \(\mathrm{X}_{i}\) 整除而 \(\mathrm{Q}_{k}\) 则否。证明传递子 J: \(\mathrm{X}_{i}\)（由满足 \(\mathrm{X}_{i} \mathrm{P} \in \mathrm{J}\) 的元素 P 所组成的 A 的理想）由单项式 \(\mathrm{X}_{i}^{-1} \mathrm{P}_{1}, \ldots, \mathrm{X}_{i}^{-1} \mathrm{P}_{r}\) 与 \(\mathrm{Q}_{1}, \ldots, \mathrm{Q}_{s}\) 生成。
\item
  由此推出：若 \(J: X_i\) 等于 \(J\) 或 \(A\) 对每个 \(i \in I\) 成立，则 \(J\) 是变量 \(X_i\)（它们属于 \(J\)）所生成的理想。
\item
  设 M 为一个 A-模，p 为一个整数。假设 \(\operatorname{Tor}_{p}^{A}(A/J',M)\)（相应地，\(\operatorname{Ext}_{A}^{p}(A/J',M)\)）对每个理想 \(J'\)（包含 J 且由某些 \(X_{i}\) 生成）均为零。证明 \(\operatorname{Tor}_{p}^{A}(A/J,M)\)（相应地，\(\operatorname{Ext}_{A}^{p}(A/J,M)\)）为零（取一个理想 \(J'\)，它包含 J、由单项式生成、使 \(\operatorname{Tor}_{p}^{A}(A/J',M) \neq 0\) 且对这些性质为极大；由正合列
\[
0 \rightarrow \mathrm{A} / (\mathrm{J} ^ {\prime}: \mathrm{X} _ {i}) \longrightarrow \mathrm{A} / \mathrm{J} ^ {\prime} \longrightarrow \mathrm{A} / (\mathrm{J} ^ {\prime} + \mathrm{AX} _ {i}) \rightarrow 0
\]

以及 b) 推出 J' 由它所包含的元素 \(X_{i}\) 生成，这与假设矛盾）。

\item
  由此推出不等式 \(\mathrm{dp}_{\mathrm{A}}(\mathrm{A}/\mathrm{J}) \leqslant \mathrm{Card}(\mathrm{I})\)。

\item
  设 A 为环，\(\mathbf{x} = (x_{i})_{i \in I}\) 为 A 中元素的一个有限族；假设对 I 的每个子集 I′，族 \((x_{i})_{i \in I'}\) 对 A 是完全割的。设 J 为由元素 \(x^{\alpha}\) 生成的 A 的理想，其中 \(\alpha\) 取遍 \(N^{1}\) 的一个有限子集 F。证明不等式 dp \(_{A}\)(A/J) ≤ Card(I)（设 J \(_{0}\) 为环 A \(_{0}\) = Z{[}(X \(_{i}\)) \(_{i \in I}\){]} 中由诸 \(X^{\alpha}\)（\(\alpha \in F\)）生成的理想；由 c) 推出 Tor \(_{p}^{A_{0}}\)(A \(_{0}\)/J \(_{0}\), A) 为零，并考虑 A \(_{0}\)-模 A \(_{0}\)/J \(_{0}\) 的一个长度 ≤ n 的投射分解来完成证明）。
\end{enumerate}

¶ 12) 设 A 为正则局部环，M 为有限生成 A-模，且 A-模 End \(_{A}\)(M) 是自由的。我们旨在证明下列条件等价：

\begin{enumerate}
\def\labelenumi{(\roman{enumi})}
\item
  A-模 M 是自由的；
\item
  A-模 M 是自反的；
\item
  \(\operatorname{Ext}_{\mathrm{A}}^{1}(\mathbf{M},\mathbf{M}) = 0\)
\end{enumerate}

\begin{enumerate}
\def\labelenumi{\alph{enumi})}
\item
  证明蕴含关系 (iii) \(\Rightarrow\) (ii)（利用 VII, § 4 习题 2）。
\item
  证明逆蕴含（当 dim(A) ≤ 2 时利用习题 4，然后利用 § 1 习题 18 对 dim(A) 作归纳论证）。
\item
  假设 \(\dim(A) \leqslant 3\)。证明蕴含关系 (ii) \(\Rightarrow\) (i)（当 \(\dim(A) \leqslant 2\) 时利用习题 4；若 \(\dim(A) = 3\)，注意 (ii) 蕴含 \(\mathrm{dp}_{A}(M) \leqslant 1\)，再注意 (iii) 与 \(n^{\circ} 3\) 的注 2 蕴含 M 是投射的）。
\item
  假设 M 是自反的。设 x 为 \(m_{A} = m_{A}^{2}\) 的一个元素；用 B 表示环 A/xA，用 N 表示 B-模 M/xM，用 \(N^{*}\) 表示它的对偶。利用 b)，证明 B-模 \(\operatorname{End}_{\mathbb{B}}(\mathbb{N})\) 是自由的，再证明 \(\operatorname{End}_{\mathbb{B}}(\mathbb{N}^{*})\) 是自由的（由 VII, § 4, 第 2 节推出 \(\operatorname{End}_{\mathbb{B}}(\mathbb{N}^{*})\) 与 \(\operatorname{End}_{\mathbb{B}}(\mathbb{N})\) 的双对偶等同）。
\item
  证明 (ii) 蕴含 (i)（对 A 的维数作归纳论证，设其 ≥ 4；采用 d) 的记号，归纳假设蕴含 B-模 N* 是自由的；由此并利用 § 1 习题 18 与第 4 节命题 7，推出 prof \(_{A}\)(Ext \(^{1}_{A}\)(M, B)) ≥ 1，再推出按归纳假设为有限长度的 Ext \(^{1}_{A}\)(M, A) 为零；断言 B-模 M*/xM* 与 N* 等同，从而是自由的，进而 A-模 M* 是自由的，最后 M 是自由的。）
\item
  由这些结果推出环 A 是因子分解整环这一事实的另一证明。
\end{enumerate}

\begin{enumerate}
\def\labelenumi{\arabic{enumi})}
\setcounter{enumi}{11}
\tightlist
\item
  设 A 为环，\((x_{1},\ldots,x_{n})\) 为 A 中元素的一个族，生成理想 I。若族 \((x_{1},\ldots,x_{n})\) 的诸 \(x_{i}\) 在 I/I \(^{2}\) 中的类构成该 A/l-模的一个基，则称该族是 A-无关的。
\end{enumerate}

\begin{enumerate}
\def\labelenumi{\alph{enumi})}
\item
  每个对 A 的完全割族都是 A-无关的。
\item
  设 \((x_{1},\ldots,x_{n})\) 为一个 A-无关族，且 \(x_{1}=x_{1}^{\prime}x_{1}^{\prime\prime}\)。证明族 \((x_{1}^{\prime},\ldots,x_{n})\) 与 \((x_{1}^{\prime\prime},\ldots,x_{n})\) 都是 A-无关的，并且有
\end{enumerate}

\[
\lg_ {\mathrm{A}} \left(\mathrm{A} / \left(x _ {1}, \dots , x _ {n}\right)\right) = \lg_ {\mathrm{A}} \left(\mathrm{A} / \left(x _ {1} ^ {\prime}, \dots , x _ {n}\right)\right) + \lg_ {\mathrm{A}} \left(\mathrm{A} / \left(x _ {1} ^ {\prime \prime}, \dots , x _ {n}\right)\right)
\]

（注意商 \((x_{1}^{\prime},\ldots,x_{n})/(x_{1},\ldots,x_{n})\) 同构于 \(A/(x_{1}^{\prime\prime},\ldots,x_{n})\)）。

¶ 13) 设 \(p\) 为素数，\(q\) 为 \(p\) 的一个幂。设 A 为特征 \(p\) 的约化 Noether 局部环。用 \(\mathrm{A}^q\) 表示由诸元素 \(a^q\)（\(a \in \mathrm{A}\)）组成的 A 的子环。证明：为使 A 正则，必须且只需 \(\mathrm{A}^q\)-模 A 是平坦的。

（化为 A 为完备的情形，于是 A 是形式幂级数代数 \(B = \kappa_{A}[[X_{1}, \ldots, X_{n}]]\) 对某个理想 b 的商，该理想可假设含于 \(m_{B}^{2}\)。若 A 正则，则 b = 0。若 A 在 \(A^{q}\) 上平坦，设 \(x_{i}\) 为 \(X_{i}\) 在 A 中的像；证明族 \((x_{1}^{q}, \ldots, x_{n}^{q})\) 是 A-无关的（习题 12），由此推出等式 \(\lg_{\mathrm{A}}(\mathrm{A}/(x_{1}^{q}, \ldots, x_{n}^{q})) = q^{n}\)。证明这蕴含 \(b \subset (X_{1}^{q}, \ldots, X_{n}^{q})\)；类似地证明对每个 s 有 \(b \subset (X_{1}^{q^{s}}, \ldots, X_{n}^{q^{s}})\)，从而最终 b = 0。）

\exercisesection{完全交}

\begin{enumerate}
\def\labelenumi{\arabic{enumi})}
\item
  设 \(k\) 为域，A 为多项式环 \(k[X,Y,Z]\) 对由 \(Z^2,ZX,Z(Y-1)\) 生成的理想的商环，\(x\) 与 \(y\) 为 X 与 Y 在 \(\Lambda\) 中的类。证明序列 \((x,y)\) 对 A 是完全割的但不是 A-正则的，而序列 \((y,x)\) 却是 A-正则的。
\item
  设 A 为正则局部环，I 为 A 的一个理想。
\end{enumerate}

\begin{enumerate}
\def\labelenumi{\alph{enumi})}
\tightlist
\item
  若 ht(I) = 1，证明下列条件等价：
\end{enumerate}

\begin{enumerate}
\def\labelenumi{(\roman{enumi})}
\item
  理想 I 是主理想；
\item
  A/I 是 Gorenstein 环；
\item
  A/I 是 Macaulay 环。
\end{enumerate}

\begin{enumerate}
\def\labelenumi{\alph{enumi})}
\setcounter{enumi}{1}
\item
  假设 \(ht(I)=2\)。为使 A/I 是 Gorenstein 环，必须且只需理想 I 是完全割的（设 \((L,p)\) 为 I 的一个极小投射分解；注意 \(L_{1}\) 的秩为 1，从而 \(L_{0}\) 的秩为 2）。
\item
  设 \(k\) 为交换域；取 \(A = k[[X,Y,Z]]\) 与 \(I = (XY,YZ,ZX, X^2 - Y^2, Y^2 - Z^2)\)。证明 \(I\) 是一个高度为 3 的不完全割的理想，而 \(A/I\) 是 Gorenstein 环。
\end{enumerate}

\begin{enumerate}
\def\labelenumi{\arabic{enumi})}
\setcounter{enumi}{2}
\tightlist
\item
  设 R 为正则局部环，I 为 R 的理想，A 为环 R/I。设 \(\mathbf{x} = (x_{1}, \ldots, x_{d})\) 为 R 的一个坐标系，\(\mathbf{y} = (y_{1}, \ldots, y_{m})\) 为 I 的一个极小生成系，于是 \(m = [I/\mathfrak{m}_{R}\mathbf{l} : \kappa_{R}]\)。
\end{enumerate}

\begin{enumerate}
\def\labelenumi{\alph{enumi})}
\item
  当 I 包含于 \(\mathfrak{m}_{\mathrm{R}}^{2}\) 时，定义 \(\mathrm{H}_1(\mathbf{x},\mathrm{A})\) 到 \(\mathrm{I} / \mathfrak{m}_{\mathrm{R}}\mathrm{I}\) 上的一个典范同构（考虑 Koszul 复形的正合列）。
\item
  为使 \(x_{i}\) 在 A 中的像构成 \(m_{A}\) 的极小生成系，必须且只需 I 包含于 \(m_{R}^{2}\)。
\item
  证明 \(\kappa_{\mathrm{A}}\)-向量空间 \(H_1(\mathbf{x}, A)\) 的维数等于 \(m - d + [\mathfrak{m}_{\mathrm{A}} / \mathfrak{m}_{\mathrm{A}}^{2} : \kappa_{\mathrm{A}}]\)。
\end{enumerate}

4) 设 A 为 Noether 局部环。令 \(\delta(A) = [\mathfrak{m}_{A}/\mathfrak{m}_{A}^{2}:\kappa_{A}] - \dim(A)\)（参见 § 4, 第 5 节）。

\begin{enumerate}
\def\labelenumi{\alph{enumi})}
\item
  证明 \(\kappa_{\Lambda}\)-向量空间 \(H_{i}(\mathbf{z}, A)\)（其中 \(\mathbf{z}\) 为 \(\mathfrak{m}_{\mathrm{A}}\) 的一个极小生成系）的维数与 \(\mathbf{z}\) 的选取无关；记之为 \(h_{i}(\mathrm{A})\)。
\item
  证明等式 \(h_i(\mathrm{A}) = h_i(\widehat{\mathrm{A}})\) 对每个 \(i\) 成立。
\item
  证明有 \(h_1(A) \geqslant \delta(A)\)，且等号成立当且仅当 \(A\) 是完全交环（第 2 节的注 2；化为 \(A\) 完备的情形，并把 \(A\) 表为商 \(R/I\)（其中 \(R\) 为正则局部环且 \(I \subset m_R^2\)）后应用习题 4）。
\item
  设 \((x_{1},\ldots,x_{p})\) 为 \(m_{A}\) 中元素的生成理想 J 的一个完全割序列。证明等式 \(h_{1}(A/J)-\delta(A/J)=h_{1}(A)-\delta(A)\)（化为 p=1 的情形，并按 \(x_{1}\) 属于或不属于 \(m_{A}^{2}\) 分两种情形讨论）。
\end{enumerate}

\begin{enumerate}
\def\labelenumi{\arabic{enumi})}
\setcounter{enumi}{4}
\tightlist
\item
  设 R 为正则局部环，I 为 R 的理想，A 为环 R/I。为使 A 是完全交环，必须且只需理想 I 是完全割的（利用习题 4 c) 与 3 c)）。
\end{enumerate}

¶ 6) 设 A 为 Noether 局部环，I 为 A 的理想。

\begin{enumerate}
\def\labelenumi{\alph{enumi})}
\tightlist
\item
  证明下列条件等价：
\end{enumerate}

\begin{enumerate}
\def\labelenumi{(\roman{enumi})}
\item
  理想 I 是完全割的；
\item
  A/I-模 I/I² 是自由的且 dp\_A(A/I) \textless{} +∞。
\end{enumerate}

在条件 (ii) 下，由 § 3 习题 2 推出存在 I 中的可消去元 \(x\)，使 \(x \notin \mathfrak{m}_{\mathrm{A}}\mathrm{I}\)；对 \(\mathrm{I} / \mathrm{I}^2\) 的维数作归纳论证。）

\begin{enumerate}
\def\labelenumi{\alph{enumi})}
\setcounter{enumi}{1}
\item
  证明 A 中使理想 \(I_{p}\) 为完全割的素理想 p 所成的集合在 Spec(A) 中是开的。
\item
  设 B 为可表现环；证明 B 中使 \(B_{q}\) 为完全交环的素理想 q 所成的集合在 Spec(B) 中是开的。
\end{enumerate}

\begin{enumerate}
\def\labelenumi{\arabic{enumi})}
\setcounter{enumi}{6}
\tightlist
\item
  设 A 为环，B 为一个 A-代数；假设 A-模 B 是有限型自由的。在以下习题中，用 B* 表示 A-模 Hom \(_{A}\)(B, A)，并赋予其自然的 B-模结构。
\end{enumerate}

\begin{enumerate}
\def\labelenumi{\alph{enumi})}
\item
  设 \((e_i)_{i\in I}\) 为 A-模 B 的一个基，\((e_i^*)_{i\in I}\) 为 \(\mathbf{B}^*\) 的对偶基。证明 \(\mathrm{Tr}_{\mathrm{B / A}}\)（\(\mathbf{B}^*\) 的元素，A, III, 第~110 页）等于 \(\sum_{j}e_{i}e_{i}^{*}\)。
\item
  假设 B-模 \(B^{*}\) 是秩 1 自由的；设 \(\ell\) 为构成该模的一个基的元素，\(\delta\) 为 B 中使 \(Tr_{B/A} = \delta\ell\) 的元素。B 的理想 \(\delta B\) 不依赖于 \(\ell\) 的选取；称之为 B 在 A 上的微差理想。证明 \(N_{B/A}(\delta)\) 是 B 在 A 上的判别式理想的一个生成元（A, III, 第 115 页；计算 B 的基（\(e_{i}\)）在 \(\Lambda\) 上的判别式，利用 \(\delta\) 在此基下的乘法矩阵）。
\end{enumerate}

\begin{enumerate}
\def\labelenumi{\arabic{enumi})}
\setcounter{enumi}{7}
\tightlist
\item
  设 A 为环，P 为次数 \(n\) 的 A{[}X{]} 首一多项式，B 为 A-代数 A{[}X{]}/(P)。用 \(x\) 表示 X 在 B 中的类。设 \(\ell: \mathrm{B} \to \mathrm{A}\) 为满足 \(\ell(x^{n-1}) = 1\)、且 \(\ell(x^i) = 0\)（\(0 \leqslant i \leqslant n-2\)）的 A-线性形式；用 \(\tilde{\ell}: \mathrm{B}[\mathrm{X}] \to \mathrm{A}[\mathrm{X}]\) 表示 A-映射{[}X{]}-线性映射（由 \(\ell\) 经纯量扩张导出）。
\end{enumerate}

\begin{enumerate}
\def\labelenumi{\alph{enumi})}
\item
  设 \(\sigma : B[X] \to A[X]\) 为 \(A[X]\)-线性映射，满足 \(\sigma(x^i) = X^i\)（\(i = 0, \ldots, n - 1\)）。证明公式 \(P\tilde{\ell}(Q) = \sigma((X - x)Q)\) 对每个元素 \(Q\)（属于 \(B[X]\)）成立（对 \(Q = 1, \ldots, x^{n-1}\) 验证之）。
\item
  设 \(\mathrm{P}^{*}(\mathrm{X}) = b_{n-1}\mathrm{X}^{n-1} + \ldots + b_{0}\) 为 B{[}X{]} 中使 \(\mathrm{P}(\mathrm{X}) = (\mathrm{X} - x)\mathrm{P}^{*}(\mathrm{X})\) 在 B{[}X{]} 中成立的多项式。证明 \((b_{0}\ell, \ldots, b_{n-1}\ell)\) 是 \(B^{*}\) 的与 B 的基 \((1, \ldots, x^{n-1})\) 对偶的基（对多项式 \(x^{i}\mathrm{P}^{*}\) 应用 a)）。由此推出 B-模 \(B^{*}\) 是秩 1 自由的，且 \(\ell\) 构成它的一个基。
\item
  证明等式 \(\mathrm{Tr}_{\mathrm{B / A}} = \mathrm{P}'(x)\ell\)（在 \(\mathrm{B}^*\) 中；参见习题 7 a)）。
\end{enumerate}

\begin{enumerate}
\def\labelenumi{\arabic{enumi})}
\setcounter{enumi}{8}
\tightlist
\item
  设 A 为环，n 为整数，S 为代数 \(A[X_{1},\ldots,X_{n}]\)（相应地，\(A[[X_{1},\ldots,X_{n}]]\)），a 为 S 的一个由对 S 完全割的序列 \(\mathbf{P}=(\mathbf{P}_{1},\ldots,\mathbf{P}_{n})\) 生成的理想。令 \(B=S/a\)，并用 \(\pi\) 表示 S 到 S/a 上的典范同态。假设 A-模 B 是有限型自由的。我们旨在证明 B-模 \(B^{*}\) 是秩 1 自由的，且 B 在 A 上的微差理想由 \(\pi(\det(\frac{\partial P_{i}}{\partial X_{j}}))\) 生成。
\end{enumerate}

\begin{enumerate}
\def\labelenumi{\alph{enumi})}
\item
  令 \(C = S \otimes_{A} B\)，并令 \(\xi_{i} = X_{i} \otimes 1 - 1 \otimes \pi(X_{i})\)（\(1 \leqslant i \leqslant n\)）。用 \(\varphi : C \to B\) 表示满足 \(\varphi(P \otimes b) = \pi(P)b\) 的 A-代数同态。证明 \(\varphi\) 的核由序列 \(\boldsymbol{\xi} = (\xi_{1}, \ldots, \xi_{n})\) 生成，且该序列对 C 是完全割的。
\item
  设 \((c_{ij}) \in \mathbf{M}_{n}(\mathbf{C})\) 为满足 \(P_{i} \otimes 1 = \sum_{j} c_{ij} \xi_{j}\) 的矩阵。证明等式 \(\pi(\frac{\partial P_{i}}{\partial X_{j}}) = \varphi(c_{ij})\)。
\item
  矩阵 \((c_{ij})\) 定义复形的一个态射 \(\mathbf{K}_{\bullet}(\mathbf{P},\mathrm{C})\to\mathbf{K}_{\bullet}(\boldsymbol{\xi},\mathrm{C})\)（A, X, 第~151 页），从而给出 S-模复形的一个态射 \(u:\mathbf{K}_{\bullet}(\mathbf{P},\mathrm{S})\to\mathbf{K}_{\bullet}(\boldsymbol{\xi},\mathrm{C})\)。同态 \(u_{n}:S\to C\) 把 \(1_{S}\) 映到 C 的元素 \(d=\det(c_{ij})\)。
\item
  设 \(v: \mathrm{Hom}_{\mathrm{S}}(\mathrm{C}, \mathrm{S}) \otimes_{\mathrm{C}} \mathrm{B} \to \mathrm{B}\) 为满足 \(v(f \otimes 1) = \pi(f(d))\) 的 B-线性同态；证明 \(v\) 是同构（把 \(v\) 与 \(\mathrm{H}^n(\mathrm{Homgr}_{\mathrm{S}}(u, 1_{\mathrm{S}}))\) 等同，并注意复形 \(\mathbf{K}_{\bullet}(\mathbf{P}, \mathrm{S})\) 与 \(\mathbf{K}_{\bullet}(\boldsymbol{\xi}, \mathrm{C})\) 中的每一个都给出 B 的一个由有限型自由 S-模组成的分解）。

\item
  考虑复合映射 \(t: \mathrm{Hom}_{\mathrm{A}}(\mathrm{B}, \mathrm{A}) \to \mathrm{Hom}_{\mathrm{S}}(\mathrm{C}, \mathrm{S}) \to \mathrm{B}\)，其中第一个箭头由纯量扩张得到，第二个由 \(v\) 得到。证明 \(t\) 是 B-线性且双射的（把 \(\mathrm{Hom}_{\mathrm{S}}(\mathrm{C}, \mathrm{S})\) 与 \(\mathrm{Hom}_{\mathrm{A}}(\mathrm{B}, \mathrm{A}) \otimes_{\mathrm{B}} \mathrm{C}\) 等同）。
\item
  证明公式 \(t(\mathrm{Tr}_{\mathrm{B / A}}) = \varphi (d)\)（用 \(d\) 与 \(\mathrm{Tr}_{\mathrm{B / A}}\) 由 B 的一个基及 \(\mathrm{B}^*\) 的对偶基表出，参见习题 7, a)）。
\item
  由 b)、e) 与 f) 推出 B 在 A 上的微差理想由 \(\pi(\det(\frac{\partial P_i}{\partial X_j}))\) 生成。
\end{enumerate}

\begin{enumerate}
\def\labelenumi{\arabic{enumi})}
\setcounter{enumi}{9}
\tightlist
\item
  设 A 为离散赋值环，v 为其赋范赋值，K 为其分式域。考虑一个 Noether A-代数 B，它带有 A-代数同态 \(\varepsilon:B\to A\)。用 I 表示 \(\varepsilon\) 的核，J 表示其在 B 中的零化子。
\end{enumerate}

\begin{enumerate}
\def\labelenumi{\alph{enumi})}
\tightlist
\item
  证明下列条件等价：
\end{enumerate}

\begin{enumerate}
\def\labelenumi{(\roman{enumi})}
\item
  \(\Lambda\)-模 \(1 / I^2\) 是有限长度的；
\item
  \(K \otimes_{A} B\) 在极大理想 \(K \otimes_{A} I\) 处的局部环（作为 K-代数）同构于 K；
\item
  我们有 \(\mathrm{K}\otimes_{\mathrm{A}}\mathrm{B} = (\mathrm{K}\otimes_{\mathrm{A}}\mathrm{I})\oplus (\mathrm{K}\otimes_{\mathrm{A}}\mathrm{J})\)
\end{enumerate}

以下假设这些条件满足；用 \(\gamma(B, \varepsilon)\)（或简作 \(\gamma(B)\)）表示 A-模 \(I/I^{2}\) 的长度，用 \(\eta(B, \varepsilon)\)（或简作 \(\eta(B)\)）表示 A-模 \(A/\varepsilon(J)\) 的长度。

\begin{enumerate}
\def\labelenumi{\alph{enumi})}
\setcounter{enumi}{1}
\item
  设 b 为 B 的一个含于 I 的理想；考虑 A-代数 \(B' = B/b\) 以及同态 \(\varepsilon': B' \to A\)（由 \(\varepsilon\) 导出）。证明 \(B'\) 也满足 a) 的条件，且有 \(\gamma(B', \varepsilon') \leqslant \gamma(B, \varepsilon)\) 与 \(\eta(B', \varepsilon') \leqslant \eta(B, \varepsilon)\)。
\item
  假设 A-模 B 是有限型自由的。证明 J 是秩 1 自由的 A-模，是 B 的一个直和因子，且 \(\mathrm{Tr}_{\mathrm{B / A}}(x) = \varepsilon (x)\) 对每个 \(x\in J\) 成立。
\item
  进而假设 B 是 Gorenstein 环；用 \(\ell\) 表示 B-模 B* 的一个基（§ 3, 习题 5），用 δ 表示 B 在 A 上的微差理想的一个生成元（习题 7, b)）。证明等式 η(B) = v(ε(δ))（证明 \(\ell\)(J) 等于 A，并利用 c)）。
\item
  置于 \(b\) 的情形，且 \(\mathfrak{b} \neq 0\)；假设 A-代数 B 与 B' 都是有限且平坦的，且 B 是 Gorenstein 环。证明有 \(\eta(B') < \eta(B)\)（注意 J 在 B/ \(\mathfrak{m}_{A}\) B 中的像是 B/ \(\mathfrak{m}_{A}\) B 的底座，故含于 \(\mathfrak{b}/\mathfrak{m}_{A}\mathfrak{b}\)）。
\end{enumerate}

¶ 11) 保留上一习题的记号；假设环 A 与 B 都是局部的且完备的。

\begin{enumerate}
\def\labelenumi{\alph{enumi})}
\item
  证明 B 同构于形式幂级数代数 \(S = A[[X_{1}, \ldots, X_{n}]]\) 对一个含于 \((X_{1}, \ldots, X_{n})\) 的理想 a 的商；可取 \(n = \dim_{\kappa_{A}}(\mathfrak{n}/\mathfrak{n}^{2})\)，其中 n 为 \(m_{B}\) 在 \(B/m_{A}B\) 中的像。
\item
  证明下列条件等价：
\end{enumerate}

\begin{enumerate}
\def\labelenumi{(\roman{enumi})}
\item
  B 是有限平坦 A-代数且是完全交环；
\item
  理想 a 由一个序列 \((\mathrm{P}_{1},\ldots,\mathrm{P}_{n})\) 生成，该序列对 \(\kappa_{A}\otimes_{A}S\) 完全割。
\end{enumerate}

（参见习题 5。）

\begin{enumerate}
\def\labelenumi{\alph{enumi})}
\setcounter{enumi}{2}
\item
  若这些条件满足，证明有 \(\eta(B) = \gamma(B) = v(\varepsilon(\det(\frac{\partial P_i}{\partial X_j})))\)（利用习题 9 与 10，\(d\)）。
\item
  证明可以找到 a 中元素的一个序列 \((f_{1},\ldots,f_{n})\)，它对 \(\kappa_{A}\otimes_{A}S\) 完全割，使 A-代数 \(\mathrm{B}^{\prime}=\mathrm{S}/(f_{1},\ldots,f_{n})\) 满足 \(\gamma(\mathrm{B}^{\prime})=\gamma(\mathrm{B})\)。
\end{enumerate}

（由正合列 \(\mathfrak{a}/\mathfrak{a}^{2}\to\Omega_{\mathrm{A}}(\mathrm{S})\otimes_{\mathrm{S}}\mathrm{B}\to\Omega_{\mathrm{A}}(\mathrm{B})\to0\)（A, III, 第~137 页）导出 A-模的正合列 \(\mathfrak{a}/\mathfrak{a}^{2}\otimes_{\mathrm{S}}\mathrm{A}\to\Omega_{\mathrm{A}}(\mathrm{S})\otimes_{\mathrm{S}}\mathrm{A}\to\mathrm{I}/\mathrm{I}^{2}\to0\)；取元素 \(f_{1},\ldots,f_{n}\)（取自 \(\mathfrak{a}\)），使其在 \(\Omega_{\mathrm{A}}(\mathrm{S})\otimes_{\mathrm{S}}\Lambda\) 中的像构成 \(\mathfrak{a}/\mathfrak{a}^{2}\otimes_{\mathrm{S}}\mathrm{A}\) 的像在 A 上的一个基。）

\begin{enumerate}
\def\labelenumi{\alph{enumi})}
\setcounter{enumi}{4}
\tightlist
\item
  于是，利用习题 10, e)，我们有 \(\gamma(B) \geqslant \eta(B)\)，且等号成立当且仅当 B 是完全交环。
\end{enumerate}

\exercisesection{正则代数中的纯量扩张}

\begin{enumerate}
\def\labelenumi{\arabic{enumi})}
\tightlist
\item
  设 \(k\) 为特征 \(p > 2\) 的非完全域，\(a\) 为 \(k - k^p\) 的一个元素，且 \(A = k[\mathrm{X},\mathrm{Y}] / (\mathrm{X}^2 -\mathrm{Y}^p +a)\)。
\end{enumerate}

\begin{enumerate}
\def\labelenumi{\alph{enumi})}
\item
  证明环 A 是正则的（可以证明 A 在其分式域 \(K = k(Y)[X]/(X^{2} - Y^{p} + a)\) 中整闭）。
\item
  证明 k 在 K 中是代数闭的。
\item
  设 \(k'\) 为 \(k(a^{1/p})\)，它是 \(k\) 的一个扩张，证明环 \(A \otimes_k k'\) 不是正规的。
\end{enumerate}

\begin{enumerate}
\def\labelenumi{\arabic{enumi})}
\setcounter{enumi}{1}
\tightlist
\item
  设 \(k\) 为域，A 为一个本质有限型 \(k\)-代数，B 为一个 \(k\)-代数，\(p\) 为整数。
\end{enumerate}

\begin{enumerate}
\def\labelenumi{\alph{enumi})}
\item
  若 A 与 B 都满足性质 \(S_{p}\)（§ 1, 习题 8），则 \(A \otimes_{k} B\) 亦然（按命题 4 的证明方式论证，并利用该处 f)）。
\item
  设 K 为 k 的一个扩张；为使 \(\mathrm{A}_{(\mathbb{K})}\) 满足 \(S_{p}\)，必须且只需 A 满足之。
\end{enumerate}

\begin{enumerate}
\def\labelenumi{\arabic{enumi})}
\setcounter{enumi}{2}
\tightlist
\item
  设 A 为 Noether 环，\(\rho : A \to B\) 为环同态。若 B 是 Noether 环且是平坦 A-模，并且对每个 \(\mathfrak{p} \in \operatorname{Spec}(A)\)，\(\kappa(\mathfrak{p}) \otimes_{A} B\) 都是绝对正则的，则称 A-代数 B 是绝对正则的。
\end{enumerate}

\begin{enumerate}
\def\labelenumi{\alph{enumi})}
\item
  设 T 为 B 的一个乘性子集，S 为 A 的一个使 \(\rho(S) \subset T\) 的子集。若 B 是正则 A-代数，则 \(T^{1}B\) 是 \(S^{-1}A\) 上的正则代数。
\item
  为使 Noether A-代数 B 是绝对正则的，必须且只需对 B 的每个极大理想 n，\(A_{\rho^{-1}(n)}\)-代数 \(B_{n}\) 都是绝对正则的。
\item
  若 C 是绝对正则 B-代数且 B 是绝对正则 A-代数，则 C 是绝对正则 A-代数。
\item
  若 B 是 Noether 环且 C 是忠实平坦 B-模又是绝对正则 A-代数，则 B 是绝对正则 A-代数。
\item
  设 B 为绝对正则 A-代数，A′ 为一个本质有限型 A-代数。则 A′-代数 A′ ⊗A B 是绝对正则的。
\end{enumerate}

\begin{enumerate}
\def\labelenumi{\arabic{enumi})}
\setcounter{enumi}{3}
\item
  设 A 为 Noether 环，B 为一个忠实平坦的绝对正则 A-代数（习题 3）。为使 A 是正则（相应地，正规、Macaulay、Gorenstein）环，必须且只需 B 亦然。
\item
  若 A 是 Noether 环，并且对 A 的每个素理想 p，局部环 \(\Lambda_{p}\) 的完备化都是绝对正则的 \(A_{p}\)-代数（习题 3），则称 A 为 Grothendieck 环。
\end{enumerate}

\begin{enumerate}
\def\labelenumi{\alph{enumi})}
\item
  证明 Grothendieck 环的每个分式环与每个商环都是 Grothendieck 环。
\item
  设 A 为 Noether 环，B 为一个忠实平坦的绝对正则 A-代数。若 B 是 Grothendieck 环，证明 A 亦然（设 p 为 A 的素理想，q 为 B 的位于 p 之上的素理想，将利用 III, § 5, 第 4 节, 命题 4 证明 \(\widehat{B}_{q}\) 是 \(\widehat{\Lambda_{p}}\)-模且忠实平坦；然后应用习题 3）。
\end{enumerate}

\begin{enumerate}
\def\labelenumi{\arabic{enumi})}
\setcounter{enumi}{5}
\item
  设 A 为 Grothendieck 环，I 为 A 的理想，\(\widehat{\mathsf{A}}\) 为 A 关于 I-进拓扑的分离完备化。证明 \(\widehat{\mathsf{A}}\) 是绝对正则 A-代数。（设 \(\mathfrak{n}\) 为 \(\widehat{\mathsf{A}}\) 的一个极大理想，\(\mathfrak{m}\) 为其在 A 中的逆像；由习题 3 推出 \(\widehat{\Lambda}_{\mathfrak{n}}\) 是绝对正则的 \(\mathrm{A}_{\mathfrak{m}}\)-代数。）
\item
  \begin{enumerate}
  \def\labelenumii{\alph{enumii})}
  \tightlist
  \item
    证明特征 0 的整 Dedekind 环是 Grothendieck 环。
  \end{enumerate}
\end{enumerate}

\begin{enumerate}
\def\labelenumi{\alph{enumi})}
\setcounter{enumi}{1}
\tightlist
\item
  设 \(k\) 为特征 \(p > 0\) 的域，且 \(k\) 在 \(k^p\) 上是无限维的，\(r\) 为整数 \(\geqslant 1\)，A 为 \(k[[\mathrm{X}_1, \ldots, \mathrm{X}_r]]\) 中由系数生成 \(k^p\) 上有限维向量空间的形式幂级数所组成的子环。证明 A 是维数 \(r\) 的正则局部环，其完备化与 \(k[[\mathrm{X}_1, \ldots, \mathrm{X}_r]]\) 等同，但它不是 Grothendieck 环（利用 III, § 3 习题 17）。
\end{enumerate}

\exercisesection{光滑代数}

\begin{enumerate}
\def\labelenumi{\arabic{enumi})}
\tightlist
\item
  设 \(k\) 为环，A 为一个线性拓扑的形式光滑 \(k\)-代数（带线性拓扑）。若 A 满足下列条件，则称 A 是形式平展（相应地，形式非分歧）的 \(k\)-代数：无论 \(k\)-代数 C（赋予离散拓扑）与 C 的平方零理想 N 如何，每个 \(k\)-代数的连续同态 A → C/N 都容许唯一（相应地，至多一个）提升 A → C。为使 A 在 \(k\) 上形式平展，必须且只需它既形式光滑又形式非分歧。\(k\) 的每个商都是形式非分歧的。
\end{enumerate}

\begin{enumerate}
\def\labelenumi{\alph{enumi})}
\item
  对形式光滑或形式平展代数证明命题 3 与命题 4 的类似陈述。
\item
  设 J 为 \(\Lambda\) 的一个理想，使 A 的拓扑是 J-进拓扑；用 \(\widehat{\Omega}_{k}(\Lambda)\) 表示 A-模 \(\Omega_{k}(A)\)（关于 J-进拓扑）的分离完备化。为使 \(k\)-拓扑代数 A 形式光滑，必须且只需 \(\widehat{\Omega}_{k}(A)\) 为零（注意这两个性质都等价于：A 到每个被 J 的某个幂零化的 A-模中的 \(k\)-导子均为零）。
\end{enumerate}

\begin{enumerate}
\def\labelenumi{\arabic{enumi})}
\setcounter{enumi}{1}
\item
  设 A 为 Noether 环，J 为 A 的理想，\(\widehat{\mathbf{A}}\) 为 A 关于 J-进拓扑的分离完备化。若 \(\Lambda\)-代数 \(\widehat{\mathbf{A}}\)（关于离散拓扑）形式光滑，则它是形式平展的（注意有 \(\Omega_{\mathbf{A}}(\widehat{\mathbf{A}}) = J^{n}\Omega_{\mathbf{A}}(\widehat{\mathbf{A}})\) 对每个整数 \(n\) 成立）。
\item
  设 K 为域 \(k\) 的一个有限生成扩张；证明下列条件等价：
\end{enumerate}

\begin{enumerate}
\def\labelenumi{(\roman{enumi})}
\item
  k-代数 K 是形式光滑的；
\item
  \(k\)-代数 K 是形式平展的；
\item
  K 是 k 的有限可分扩张。
\end{enumerate}

\begin{enumerate}
\def\labelenumi{\arabic{enumi})}
\setcounter{enumi}{3}
\tightlist
\item
  设 \(k\) 为 Noether 环，A 为一个有限生成的 \(k\)-代数。若 \(\Lambda\) 赋予离散拓扑后是形式平展（相应地，形式光滑）的，则称它是 \(k\)-平展（相应地，光滑）代数。
\end{enumerate}

\begin{enumerate}
\def\labelenumi{\alph{enumi})}
\tightlist
\item
  证明下列条件等价：
\end{enumerate}

\begin{enumerate}
\def\labelenumi{(\roman{enumi})}
\item
  k-代数 A 是光滑的；
\item
  \(\Omega_k(A) = 0\)
\item
  \(\mathrm{A}\otimes_{k}\mathrm{A}\)-模 A 是投射的。
\end{enumerate}

（设 I 为典范同态 \(\mathrm{A} \otimes_k \mathrm{A} \to \mathrm{A}\) 的核；证明 (ii) 与 (iii) 都等价于

\begin{enumerate}
\def\labelenumi{(\roman{enumi})}
\setcounter{enumi}{3}
\tightlist
\item
  \(I_{q} = 0\) 对 A \(\otimes_{k}\) A 的每个包含 I 的素理想 q 成立。）
\end{enumerate}

\begin{enumerate}
\def\labelenumi{\alph{enumi})}
\setcounter{enumi}{1}
\item
  设 \(f \in k[T]\) 为一个使 f 与 \(f'\) 生成 k{[}T{]} 的单位理想的多项式。证明 k-代数 \(k[T]/(f)\) 是平展的（参见第 3 节，例 3）。
\item
  若 k 是域，则光滑 k-代数就是平展 k-代数，且此概念与 A, V, 第~28 页引入的概念一致（k 上有限次代数的情形由该处第~32 页定理 3 得到；若 A 是光滑 k-代数，考虑 A 的一个极大理想 m，并对 A/m \(^{n}\) 应用前一情形）。于是光滑 k-代数就是 k 的有限可分扩张的有限积（A, V, 第~34 页，定理 4）。
\item
  为使 k-代数 A 是光滑的，必须且只需对 k 的每个素理想 p，\(\kappa(\mathfrak{p})\)-代数 \(\kappa(\mathfrak{p})\otimes_{k}A\) 都是平展的。
\end{enumerate}

\begin{enumerate}
\def\labelenumi{\arabic{enumi})}
\setcounter{enumi}{4}
\tightlist
\item
  设 A 为完备 Noether 局部环，B 为一个本质有限型 A-代数。
\end{enumerate}

\begin{enumerate}
\def\labelenumi{\alph{enumi})}
\item
  证明 A 中使环 \(\Lambda_{\mathfrak{p}}\) 正则的素理想 p 所成的集合在 \(\operatorname{Spec}(\mathbf{A})\) 中是开的（将利用 VIII, § 5 习题 16 化为 A 是整环的情形，再利用 IX, § 2, 第 5 节, 定理 3 与 VIII, § 5 习题 17）。
\item
  假设 A 包含一个域 \(^{1}\)。证明 B 中使环 \(B_{q}\) 正则的素理想 q 所成的集合在 Spec(B) 中是开的（将利用 § 7, 第 9 节, 定理 3 的系 2 与 IX, § 3, 第 3 节, 定理 2）。
\end{enumerate}

\begin{enumerate}
\def\labelenumi{\arabic{enumi})}
\setcounter{enumi}{5}
\item
  设 \(k\) 为环，A 为一个线性拓扑的形式光滑 \(k\)-代数。证明对 A 的每个开理想 J，A/J-模 \((\mathrm{A} / \mathrm{J}) \otimes_{\mathrm{A}} \Omega_k(\mathrm{A})\) 是投射的（利用第 1 节的例子证明：对 A/J-模的每个满同态 \(\mathrm{M} \to \mathrm{M}''\)，每个 \(k\)-导子（自 A 到 \(\mathrm{M}''\)）都可提升到 M）。
\item
  设 \(k\) 为环，\(\rho : A \to B\) 为 \(k\)-代数同态，\(J\) 为 \(B\) 的理想。称 \(B\)（赋予 \(J\)-进拓扑）在 \(A\) 上关于 \(k\) 形式光滑，如果对每个 \(A\)-代数 \(C\)（赋予离散拓扑）与每个平方为零的理想 \(N\)（在 \(C\) 中），每个 \(A\)-代数的连续同态 \(B \to C/N\)，只要它可提升为一个 \(k\)-代数连续同态（从 \(B\) 到 \(C\) 中），就也可提升为一个 \(A\)-代数连续同态（从 \(B\) 到 \(C\) 中）。
\end{enumerate}

\begin{enumerate}
\def\labelenumi{\alph{enumi})}
\tightlist
\item
  证明下列条件等价：
\end{enumerate}

\begin{enumerate}
\def\labelenumi{(\roman{enumi})}
\item
  B 关于 J-进拓扑在 A 上关于 k 形式光滑。
\item
  对每个 \(k\)-导子 D（从 \(\Lambda\) 到一个被 J 的某个幂零化的 B-模 E 中），存在 \(k\)-导子 \(\tilde{D}: B \to E\) 使 \(\tilde{D} \circ \rho = D\)；
\item
  对每个整数 \(n \geqslant 0\)，典范映射 \(u_n: \Omega_k(A) \otimes_A (B/J^n) \longrightarrow \Omega_k(B) \otimes_B (B/J^n)\) 是单射且其像是直和项。
\end{enumerate}

\begin{enumerate}
\def\labelenumi{\alph{enumi})}
\setcounter{enumi}{1}
\tightlist
\item
  假设 B 关于 J-进拓扑在 k 上形式光滑。证明：为使 B 在 A 上形式光滑，必须且只需典范同态 \(u_{1}:\Omega_{k}(A)\otimes_{A}(B/J)\to\Omega_{k}(B)\otimes_{B}(B/J)\) 是单射且其像是直和项（利用 a) 与习题 6 证明 \(u_{1}\) 的一个收缩可提升为每个 \(u_{n}\) 的收缩）。
\end{enumerate}

¶ 8) 设 k 为域，T 为一个有限的未定元族。给 k-代数 \(k[[T]]\) 赋予离散拓扑。

\begin{enumerate}
\def\labelenumi{\alph{enumi})}
\item
  证明 \(k[[T]]\) 在 \(k[T]\) 上关于 k 形式光滑（习题 7）。
\item
  若 \(k[[\mathbf{T}]]\) 在 \(k\) 上形式光滑，证明特征 \(p\)（\(k\) 的）严格为正且 \(k\) 是 \(k^p\) 的有限扩张（由习题 2 推出 \(\Omega_{k|\mathbf{T}|}(k[[\mathbf{T}]])\) 为零，从而 \(\Omega_{k(\mathbf{T})}(k((\mathbf{T})))\) 为零；然后利用 A, V, 第~99 页命题 6 证明：\(k((\mathbf{T}))\) 中由系数生成 \(k^p\) 上有限生成扩张的幂级数组成的中介扩张等于 \(k((\mathbf{T})))\)。
\end{enumerate}

\begin{enumerate}
\def\labelenumi{\arabic{enumi})}
\setcounter{enumi}{8}
\tightlist
\item
  设 \(k\) 为特征 \(p > 0\) 的域，包含子域的一个递减有向族 \((k_{\lambda})_{\lambda \in \Lambda}\)，满足 \(\bigcap_{\lambda \in \Lambda} k_{\lambda} = k^{p}\)。
\end{enumerate}

\begin{enumerate}
\def\labelenumi{\alph{enumi})}
\item
  证明 \(\Omega(k)\) 到 \(\Omega_{k_{\lambda}}(k)\) 中的典范同态（\(\lambda\) 取遍 \(\Lambda\)）的核的交化为 (0)（考虑一个 \(p\)-基，取自 \(k\)）。
\item
  设 A 为正则局部 k-代数，赋予 \(m_{A}\)-进拓扑。证明：为使 A 在 k 上形式光滑，必须且只需 \(\Lambda\) 在 k 上关于 \(k_{\lambda}\) 形式光滑对每个 \(\lambda \in \Lambda\) 成立（利用 a)、命题 8 与习题 7, a)）。
\end{enumerate}

¶ 10) 设 A 为完备 Noether 局部环，K 为其分式域。设 p 为 A 的素理想；用 B 表示局部环 \(A_{p}\)，\(\widehat{B}\) 表示其完备化。

\begin{enumerate}
\def\labelenumi{\alph{enumi})}
\item
  假设 A 正则且 K 的特征为零。证明 K-代数 \(K \otimes_{B} \widehat{B}\) 是绝对正则的。
\item
  假设 A 正则且 K 的特征 p \textgreater{} 0，于是 A 与形式幂级数环 \(k[[T_{1}, \ldots, T_{n}]]\) 等同（IX, § 3, 第 3 节, 定理 2）。用 \((k_{\lambda})_{\lambda \subset \Lambda}\) 表示子域 \(k_{\lambda}\)（k 的子域，包含 \(k^{p}\)，且使 k 是 \(k_{\lambda}\) 的有限次扩张）所成的族。设 \(\lambda \in \Lambda\)；用 \(A_{\lambda}\) 表示环 \(k_{\lambda}[[T_{1}^{p}, \ldots, T_{n}^{p}]]\)，\(B_{\lambda}\) 表示 \(A_{\lambda}\) 在素理想 \(p \cap A_{\lambda}\) 处的局部环，\(K_{\lambda}\) 表示 \(\Lambda_{\lambda}\) 与 \(B_{\lambda}\) 的分式域。证明 A-代数 B（赋予离散拓扑）关于 \(B_{\lambda}\) 形式光滑（可先证明 \(\widehat{B}\) 同构于 \(B \otimes_{B_{\lambda}} \widehat{B_{\lambda}}\)）。
\item
  在 b) 的假设下，证明 \(K \otimes_{B} \widehat{B}\) 的每个局部环都形如 \((\widehat{B})_{\mathbf{r}}\)，其中 \(\mathfrak{r} \in \operatorname{Spec}(\widehat{B})\) 且 \(\mathfrak{r} \cap B = \{0\}\)，并且它在 K 上关于 \(K_{\lambda}\) 形式光滑。由此推出 K-代数 \(K \otimes_{B} \widehat{B}\) 是绝对正则的（证明 K 的子域族 \(\{K_{\lambda}\}_{\lambda \in \Lambda}\) 是递减有向的，且 \(\bigcap_{\lambda \in \Lambda} K_{\lambda} = k^{p}((T_{1}, \ldots, T_{n}))\)；然后应用习题 9）。
\item
  回到一般情形。证明 A 是 Grothendieck 环（§ 6, 习题 5）；若 \(\mathfrak{q} \in \operatorname{Spec}(B)\)，需要证明 \(\kappa(\mathfrak{q})\)-代数 \(\kappa(\mathfrak{q}) \otimes_{\mathrm{B}} \widehat{\mathrm{B}}\) 是绝对正则的；把 A 换成 \(\mathrm{A}/(\mathrm{A} \cap \mathfrak{q})\) 后可假设 \(\Lambda\) 是整环且 \(\mathfrak{q} = \{0\}\)；将利用 IX, § 2, 第 3 节, 定理 3 与 § 3, 第 3 节, 定理 2 化为 A 是完备正则局部环的情形。
\end{enumerate}

\begin{enumerate}
\def\labelenumi{\arabic{enumi})}
\setcounter{enumi}{10}
\item
  设 A 为 Noether 环；假设对 A 的每个极大理想 m，局部环 \(A_{m}\) 的完备化都是绝对正则的 \(A_{m}\)-代数。证明 A 是 Grothendieck 环（§ 6, 习题 5）。（利用该处与前一习题证明：对 A 的每个极大理想 m，环 \(A_{m}\) 都是 Grothendieck 环。）
\item
  设 A 为 Noether 局部环。为使 A 是 Grothendieck 环，必须且只需：对每个有限整 A-代数 B、B 的每个极大理想 m 以及完备化环 \(\widehat{B_{m}}\) 的每个满足 \(B_{m} \cap q = (0)\) 的素理想 q，局部环 \((\widehat{B_{m}})_{q}\) 都是正则的（为证此条件必要，由前一习题只需证明：对 A 的每个素理想 p 与 \(\kappa(\mathfrak{p})\) 的每个有限扩张 K，环 \(K \otimes_{A} \widehat{A}\) 都是正则的；可引入一个以 \(K\) 为分式域的有限整 A-代数 B，记 \(\mathfrak{m}_{1},\ldots,\mathfrak{m}_{r}\) 为其极大理想，并把 \(K\otimes_{A}\widehat{A}\) 的局部环解释为某个 \(\widehat{B_{\mathfrak{m}_{i}}}\) 的局部环。
\item
  设 \(k\) 为域，A 为一个本质有限型 \(k\)-代数。证明 A 是 Grothendieck 环（化为 A 是 \(k\) 上有限个未定元的多项式环的局部环的情形；利用前一习题，再化为证明：对每个素理想 \(\mathfrak{r}\)（属于 \(\mathrm{A}[\mathrm{T}_1,\ldots ,\mathrm{T}_n]\)）、\(\mathrm{A}[\mathrm{T}_1,\ldots ,\mathrm{T}_n]\) 在包含 \(\mathfrak{r}\) 的极大理想处的每个局部环 C，以及每个理想 \(\mathfrak{q}\)（属于完备化环 \(\widehat{\mathrm{C}}\)，满足 \(\mathrm{C}\cap \mathfrak{q} = \mathfrak{r}\mathrm{C}\)），环 \((\widehat{\mathrm{C}})_q / \mathfrak{r}(\widehat{\mathrm{C}})_q\) 都是正则的。为此可利用第 9 节定理 3）。
\end{enumerate}

¶ 14) 设 A 为局部 Grothendieck 环。证明 A 中使环 \(A_{p}\) 正则的素理想 p 所成的集合在 \(\operatorname{Spec}(A)\) 中是开的（设 \(\widehat{A}\) 为 A 的完备化；将利用习题 5、VIII, § 5 习题 16，以及典范映射 \(\operatorname{Spec}(\widehat{A}) \to \operatorname{Spec}(A)\) 是开的这一事实）。

¶ 15) 设 \(k\) 为特征 \(p > 0\) 的完全域，\((\mathrm{T}_{n})_{n \in \mathbf{N}}\) 为一个未定元族，K 为域 \(k((\mathrm{T}_{n})_{n \in \mathbf{N}})\)。设 \(\Lambda\) 为形式幂级数环 \(\mathrm{K}[[\mathrm{X},\mathrm{Y}]]\) 的由 \(\mathrm{K}^{p}[[\mathrm{X},\mathrm{Y}]]\) 与 K 生成的子环。环 A 是维数 2 的 Noether 正则局部环，其完备化为 \(\mathrm{K}[[\mathrm{X},\mathrm{Y}]]\)。对每个整数 \(n\)，用 \(p_{n}\) 表示 A 的元素 \(\mathrm{X} - \mathrm{T}_{n}\mathrm{Y}\)，\(q_{n}\) 表示乘积 \(p_{0} \ldots p_{n}\)，并令 \(c = \sum_{n \in \mathbf{N}} q_{n}\mathrm{T}_{n}\)。用 B 表示 \(\mathrm{K}[[\mathrm{X},\mathrm{Y}]]\) 的由 \(c\) 生成的子 A-代数。

\begin{enumerate}
\def\labelenumi{\alph{enumi})}
\item
  设 \(n\) 为整数 \(> 0\)，\(\mathfrak{p}\) 为 \(B\) 的一个包含 \(p_n\) 的高度 1 的素理想。证明环 \(B_{\mathfrak{p}}\) 不是正规的（对每个整数 \(m\)，令 \(c_m = (c - \sum_{i=0}^{m-1} q_i T_i)/q_m\)，将证明 \(B_{\mathfrak{p}} = A_{A p_n}[c_{n-1}]\)，于是 \(c_n\) 不属于 \(B_{\mathfrak{p}}\)，而 \(c_n^p\) 属于 \(A\)）。
\item
  由此推出：B 中使环 \(B_{p}\) 正则的素理想 p 所成的集合不包含 Spec(B) 的任何非空开子集（注意 B 是忠实平坦 A-模）。
\end{enumerate}

\exercisesection{有限长度模的对偶性}

\begin{enumerate}
\def\labelenumi{\arabic{enumi})}
\tightlist
\item
  设 A 为 Noether 局部环，M 为一个 A-模。证明：为使 M 是 Artin 的，必须且只需其底座 S 在 \(\kappa_{\Lambda}\) 上是有限维的，且 M 是 S 的本性扩张（也就是说，据（A, II, 第~185 页, 习题 15），M 的每个非零子模都含有 S 的一个非零元素）。
\end{enumerate}

¶ 2) 设 A 为完备 Noether 局部环，M 为一个 Λ-模。证明下列条件等价：

\begin{enumerate}
\def\labelenumi{(\roman{enumi})}
\item
  典范同态 \(\alpha_{\mathrm{M}}: \mathrm{M} \to \mathrm{D}(\mathrm{D}(\mathrm{M}))\) 是双射的；
\item
  存在正合列 \(0 \to M' \to M \to M'' \to 0\)，其中 \(M'\) 是有限型的而 \(M''\) 是 Artin 的。
\end{enumerate}

（若 M 满足 (i)，证明其底座在 \(\kappa_{A}\) 上是有限维的，且 M 的每个商亦然；构造 M 的一个有限生成子模 N，使 M/mN 是其底座的本性扩张，并利用习题 1）。

\begin{enumerate}
\def\labelenumi{\arabic{enumi})}
\setcounter{enumi}{2}
\tightlist
\item
  设 \(k\) 为域，S 为环 \(k[\mathrm{T}_{1},\ldots ,\mathrm{T}_{d}]\)，\(\mathfrak{m}\) 为理想 \((\mathrm{T}_1,\dots ,\mathrm{T}_d)\)。用 \(\mathrm{S^{*gr}}\) 表示 S 的分次对偶（第 6 节），\((\mathbf{T}^{-\alpha})_{\alpha \in \mathbf{N}^d}\) 表示 \((\mathbf{T}^{\alpha})_{\alpha \in \mathbf{N}^d}\) 的对偶基（在该处记作 \((u_{\alpha})_{\alpha \in \mathbf{N}^d}\)）。
\end{enumerate}

\begin{enumerate}
\def\labelenumi{\alph{enumi})}
\item
  对 S*gr 的每个有限生成子 S-模 M，对应 S 的理想 Ann(M)。证明这在 S*gr 的有限生成子模与 S 的含于 m 且使 S/a 为有限长度的理想 a 之间定义了一个双射对应；逆双射把理想 a 对应到 S*gr 中由被 a 零化的元素组成的子模 M（注意 M 是一个 Matlis S/a-模）。
\item
  为使 S/a 是 Gorenstein 环，必须且只需 S-模 M 是循环的。
\item
  设 \(f\) 为元素 \(\sum \mathrm{T}_i^{-2}\)（属于 \(\mathrm{S}^{*\mathrm{gr}}\)）；证明理想 \(\mathfrak{a}\)（它对应于子模 \(A f\)，后者属于 \(\mathrm{S}^{*\mathrm{gr}}\)）由二次型 \(\mathrm{T}_i\mathrm{T}_j\)（对 \(i < j\)）与 \(\mathrm{T}_i^2 - \mathrm{T}_1^2\)（对 \(i > 1\)）生成。若 \(d \geqslant 3\)，理想 \(\mathfrak{a}\) 不是完全割的；Gorenstein 环 \(A / \mathfrak{a}\) 不是完全交环（参见 § 5, 习题 5）。
\end{enumerate}

\begin{enumerate}
\def\labelenumi{\arabic{enumi})}
\setcounter{enumi}{3}
\item
  设 A 为 Noether 环，M 为有限生成 A-模。对每个 \(\mathfrak{p} \in \operatorname{Spec}(A)\)，令 \(\boldsymbol{\mu}^{i}(\mathfrak{p}, M) = \dim_{\kappa(\mathfrak{p})} \operatorname{Ext}_{A_{\mathfrak{p}}}^{i}(\kappa(\mathfrak{p}), M_{\mathfrak{p}})\)；用 \((I(\mathfrak{p}), e_{\mathfrak{p}})\) 表示 A-模 \(A/\mathfrak{p}\) 的一个内射包。设 I 为 M 的一个极小内射分解（即 \(I^{n}\) 都是 \(Z^{n}(I)\) 的内射包对每个 \(n \geqslant 0\) 成立，参见 A, X, 第 182 页, 习题 13）。证明 A-模 \(I^{p}\) 同构于 \(\bigoplus_{\mathfrak{p}} I(\mathfrak{p})^{\mu^{p}(\mathfrak{p}, M)}\)。
\item
  设 A 为 Noether 局部环，I 为一个 Matlis A-模。对每个 A-模 M，用 D(M) 表示 Matlis 对偶 \(\mathrm{Hom}_{\mathrm{A}}(\mathrm{M},\mathrm{I})\)。回顾用 \(\mathrm{dpl}_{\mathrm{A}}(\mathrm{M})\) 表示 M 的平坦分解的长度的下确界（A, X, 第 202 页, 习题 7）。证明等式 \(\mathrm{di}_{\mathrm{A}}(\mathrm{D}(\mathrm{M})) = \mathrm{dpl}_{\mathrm{A}}(\mathrm{M})\)，\(\mathrm{di}_{\mathrm{A}}(\mathrm{M}) = \mathrm{dpl}_{\mathrm{A}}(\mathrm{D}(\mathrm{M}))\)。
\item
  设 A 为 Noether 环，I 与 J 为内射 A-模。证明 A-模 \(\mathrm{Hom}_{\mathrm{A}}(\mathrm{I},\mathrm{J})\) 是平坦的（参见命题 6, b)）。
\item
  设 A 与 B 为不必交换的环，J 为一个 (A, B)-双模，P 为一个平坦的左 B-模。假设环 A 是左 Noether 的。证明：若 J 是内射 A-模，则 J ⊗B P 亦然。
\end{enumerate}

¶ 8) 设 A 为 Noether 局部环。

\begin{enumerate}
\def\labelenumi{\alph{enumi})}
\tightlist
\item
  设 M 与 N 为满足 \(\mathrm{dpl}_{\mathrm{A}}(\mathrm{M}) < +\infty\) 与 \(\operatorname{Tor}_{i}^{\mathrm{A}}(\mathrm{M}, \mathrm{N}) = 0\)（对 i \textgreater{} 0）的 A-模。对每个整数 n 定义一个典范同构
\end{enumerate}

\[
\operatorname{Ext} _ {\mathrm{A}} ^ {n} \left(\kappa_ {\mathrm{A}}, \mathrm{M} \otimes_ {\mathrm{A}} \mathrm{N}\right) \longrightarrow \bigoplus_ {p \geqslant n} \left(\operatorname{Tor} _ {p - n} ^ {\mathrm{A}} \left(\kappa_ {\Lambda}, \mathrm{M}\right) \otimes_ {k} \operatorname{Ext} _ {\mathrm{A}} ^ {p} \left(\kappa_ {\mathrm{A}}, \mathrm{N}\right)\right),
\]

并证明有 \(\mathrm{di}_{\mathrm{A}}(\mathrm{M}\otimes_{\mathrm{A}}\mathrm{N})\leqslant \mathrm{di}_{\mathrm{A}}(\mathrm{N})\)（借助习题 7 仿照 § 3 习题 6 的证明）。

\begin{enumerate}
\def\labelenumi{\alph{enumi})}
\setcounter{enumi}{1}
\tightlist
\item
  设 N, P 为满足 \(\mathrm{di}_{\mathrm{A}}(\mathrm{P}) < +\infty\) 与 \(\mathrm{Ext}_{\mathrm{A}}^{i}(\mathrm{N},\mathrm{P}) = 0\)（对 i \textgreater{} 0）的 A-模。对每个整数 n 定义一个典范同构
\end{enumerate}

\[
\operatorname{Tor} _ {n} ^ {\mathrm{A}} \left(\kappa_ {\mathrm{A}}, \operatorname{Hom} _ {\mathrm{A}} (\mathrm{N}, \mathrm{P})\right) \longrightarrow \bigoplus_ {p \geqslant n} \operatorname{Hom} _ {\kappa_ {\mathrm{A}}} \left(\operatorname{Ext} _ {\mathrm{A}} ^ {p - n} \left(\kappa_ {\mathrm{A}}, \mathrm{P}\right), \operatorname{Ext} _ {\mathrm{A}} ^ {p} \left(\kappa_ {\mathrm{A}}, \mathrm{N}\right)\right)
\]

（应用 Matlis 对偶）。

\begin{enumerate}
\def\labelenumi{\alph{enumi})}
\setcounter{enumi}{2}
\item
  证明：若存在一个内射维数与平坦维数都有限的 A-模 M，则 A 是 Gorenstein 环（在 a) 中取 N = A）。
\item
  反之，若 A 是 Gorenstein 环，证明内射维数有限的模与平坦维数有限的模是一致的（借助 a) 证明平坦维数有限的模其内射维数有限，并由 Matlis 对偶得到其逆）。
\item
  由 d) 推出下列条件等价：
\end{enumerate}

\begin{enumerate}
\def\labelenumi{(\roman{enumi})}
\item
  A 是 Gorenstein 环；
\item
  内射维数有限的有限生成 A-模与投射维数有限的有限生成 A-模相一致；
\item
  存在一个内射维数与投射维数都有限的有限生成 A-模。
\end{enumerate}

\begin{enumerate}
\def\labelenumi{\arabic{enumi})}
\setcounter{enumi}{8}
\tightlist
\item
  设 A 为满足条件 (GM) 的 Noether 局部环（§ 3, 习题 13）。
\end{enumerate}

\begin{enumerate}
\def\labelenumi{\alph{enumi})}
\item
  设 I 为内射 A-模的一个长度为 \(\ell\) 的复形，其同调非零且为有限型的。证明不等式 \(\ell \geqslant \dim(A)\)（用 Matlis 对偶化为该处 a) 的情形）。
\item
  若存在一个内射维数有限的有限生成 A-模 M，则 A 是 Macaulay 环（对 M 的一个内射分解应用 a)）。
\item
  为使 A 是 Gorenstein 环，必须且只需它拥有一个内射维数有限的循环模（利用习题 8, b)）。
\end{enumerate}

\begin{enumerate}
\def\labelenumi{\arabic{enumi})}
\setcounter{enumi}{9}
\tightlist
\item
  设 A 为 Noether 环，M, P, T 为 A-模；假设 M 是有限生成的，T 的内射维数有限，且模 \(\operatorname{Ext}_{\mathrm{A}}^{i}(\mathrm{M},\mathrm{P})\) 与 \(\operatorname{Ext}_{\mathrm{A}}^{i}(\mathrm{P},\mathrm{T})\) 对 i \textgreater{} 0 均为零。证明典范同态 \(\mathrm{M}\otimes_{\mathrm{A}}\mathrm{Hom}_{\mathrm{A}}(\mathrm{P},\mathrm{T})\to\mathrm{Hom}_{\mathrm{A}}(\mathrm{Hom}_{\mathrm{A}}(\mathrm{M},\mathrm{P}),\mathrm{T})\)（它把 \(m\otimes h\) 映为同态 \(u\mapsto h(u(m))\)）是同构（按命题 6 的证明方式论证，把 J 换成 T 的一个内射分解）。
\end{enumerate}

\exercisesection{对偶化模}

\begin{enumerate}
\def\labelenumi{\arabic{enumi})}
\tightlist
\item
  设 A 为局部 Macaulay 环。
\end{enumerate}

\begin{enumerate}
\def\labelenumi{\alph{enumi})}
\item
  假设 A 具有对偶化模 \(\Omega\)；赋予 A-模 \(\mathrm{A} \oplus \Omega\) 一个 A-代数结构，使 \(\Omega\) 成为平方为零的理想。证明 \(\mathrm{A} \oplus \Omega\) 是 Gorenstein 环（利用极大正则序列化为 A 为 Artin 环的情形；注意 \(\mathrm{A} \oplus \Omega\) 的底座与 A-模 \(\Omega\) 的底座等同）。
\item
  由此推出，为使 A 具有对偶化模，必须且只需存在一个 Gorenstein 环 B 以及从 B 到 A 的满同态。
\end{enumerate}

\begin{enumerate}
\def\labelenumi{\arabic{enumi})}
\setcounter{enumi}{1}
\tightlist
\item
  设 A 为正则局部环，I 为 A 的理想，B 为环 A/I；假设 B 是 Macaulay 环。设 (L,p) 为 A-模 B 的极小投射分解；其长度为 \(c = \text{ht}(I)\)。
\end{enumerate}

\begin{enumerate}
\def\labelenumi{\alph{enumi})}
\item
  设 \(\mathrm{L}^*\) 为复形 \(\operatorname{Homgr}_{\mathrm{A}}(\mathrm{L},\mathrm{A})\)；证明有 \(\mathrm{H}^{i}(\mathrm{L}^{*}) = 0\) 对 \(i\neq c\) 成立，且 B-模 \(\mathrm{H}^c (\mathrm{L}^*)\) 是对偶化的。
\item
  证明下列条件等价：
\end{enumerate}

\begin{enumerate}
\def\labelenumi{(\roman{enumi})}
\item
  A 是 Gorenstein 环；
\item
  复形 L 与 \(L^{*}(-d)\) 同构；
\item
  \(L_{c}\) 是秩为 1 的自由模。
\end{enumerate}

\begin{enumerate}
\def\labelenumi{\arabic{enumi})}
\setcounter{enumi}{2}
\tightlist
\item
  设 A 为具有对偶化模 \(\Omega\) 的环。
\end{enumerate}

\begin{enumerate}
\def\labelenumi{\alph{enumi})}
\item
  为使 \(\Omega\) 同构于 A 的一个理想 I，必须且只需 \(A_{p}\) 是 Gorenstein 环对 A 的每个极小素理想 \(p\) 成立（注意 \(\Omega\) 同构于 \(\bigoplus \Omega_{p}\) 的一个子模）。
\item
  当这些条件满足时，证明 I 等于 A 或者高度为 1，且 A/I 是 Gorenstein 环（先计算 A/I 的深度，再计算 A/I-模 \(\mathrm{Ext}_{\mathrm{A}}^{1}(\mathrm{A / I,I})\)）。
\item
  证明具有对偶化模的因子分解整环是 Gorenstein 环。
\end{enumerate}

\begin{enumerate}
\def\labelenumi{\arabic{enumi})}
\setcounter{enumi}{3}
\item
  设 A 为维数 \(d\) 的局部 Macaulay 环；设 \(\Omega\) 为一个有限生成的 A-模，其内射维数有限、深度为 \(d\)，且使典范同态 \(\mathrm{A} \to \mathrm{End}_{\mathrm{A}}(\Omega)\) 是双射。证明 \(\Lambda\)-模 \(\Omega\) 是对偶化的（对 \(d\) 作归纳，并由 § 3 习题 7 推出：对每个 \(\Omega\)-正则的元素 \(x\)（属于 \(\mathfrak{m}_{\mathrm{A}}\)），典范同态 \(\mathrm{A}/x\mathrm{A} \to \mathrm{End}_{\mathrm{A}/x\mathrm{A}}(\Omega/x\Omega)\) 是双射）。
\item
  设 A 为具有对偶化模 \(\Omega\) 的 Noether 局部环。设 T 为内射维数有限的有限生成 A-模。
\end{enumerate}

\begin{enumerate}
\def\labelenumi{\alph{enumi})}
\item
  证明 \(\mathrm{Ext}_{\mathrm{A}}^{i}(\Omega, \mathrm{T})\) 为零对 \(i > 0\) 成立（利用 § 3 习题 7）。
\item
  证明典范同态 \(\Omega\otimes_{\Lambda}\mathrm{Hom}_{\mathbb{A}}(\Omega,\mathrm{T})\to\mathrm{T}\) 是双射（应用 a) 与 § 8 习题 10）。
\item
  试利用 § 8 习题 8, b) 证明等式 dp \(_{A}\) Hom \(_{A}\) (Ω,T)) = dim(A) - prof(T) 与 prof(Hom \(_{A}\) (Ω,T)) = prof(T)。由此推出，若 prof(T) = dim(A)，则 A-模 T 是若干同构于 Ω 的模的直和。
\end{enumerate}

\(d)\) 设 \(c = \dim(A) - \operatorname{prof}(T)\)。证明存在整数 \(n_0, \ldots, n_c\) 以及一个正合序列

\[
0 \rightarrow \Omega^ {n _ {c}} \rightarrow \dots \rightarrow \Omega^ {n _ {0}} \rightarrow \mathrm{T} \rightarrow 0
\]

（对 c 作归纳，借助 b) 构造一个满同态 \(\Omega^{n_{0}} \to T\)）。

¶ 6) 设 A 为具有对偶化模 Ω 的局部 Noether 环；设 M 为投射维数有限的有限生成 A-模。

\begin{enumerate}
\def\labelenumi{\alph{enumi})}
\item
  证明有 \(\operatorname{Tor}_{i}^{\mathrm{A}}(\Omega, \mathrm{M}) = 0\) 对 \(i > 0\) 成立（对 \(\dim(\mathrm{A})\) 作归纳，考察一个正合序列 \(0 \to \mathrm{N} \to \mathrm{L} \to \mathrm{M} \to 0\)，其中 \(\mathrm{L}\) 是有限型自由模。若 \(x\) 是 \(\mathfrak{m}_{\mathrm{A}}\) 中一个非零因子的元素，则由归纳假设推出 \(\operatorname{Tor}_{i}^{\mathrm{A}}(\Omega, \mathrm{N})\) 当 \(i > 0\) 时为零，由此得 \(\operatorname{Tor}_{i}^{\mathrm{A}}(\Omega, \mathrm{M}) = 0\) 对 \(i > 1\) 成立；又比率为 \(x\) 的乘法映射在 \(\Omega \otimes_{\mathrm{A}} \mathrm{N}\) 中是单射，由此推出每个素理想 \(\mathfrak{p}\)（与 \(\operatorname{Tor}_{1}^{\mathrm{A}}(\Omega, \mathrm{M})\) 相伴）都是 \(\operatorname{Spec}(\mathrm{A})\) 中的极小理想；注意在这样的 \(\mathfrak{p}\) 处，\(\mathrm{A}_{\mathfrak{p}}\)-模 \(\mathrm{M}_{\mathfrak{p}}\) 是自由的）。
\item
  证明 A-模 \(\Omega\otimes_{\mathrm{A}}\mathrm{M}\) 的内射维数有限。
\item
  证明映射 \(\theta: \mathrm{M} \to \mathrm{Hom}_{\Lambda}(\Omega, \Omega \otimes_{\Lambda} \mathrm{M})\)（它由 \(\theta(m)(w) = w \otimes m\) 对 \(m \in \mathrm{M}\)、\(w \in \Omega\) 定义）是同构（借助一个投射分解化为 M 为自由模的情形）。
\end{enumerate}

\(d\) ) 映射 \([\mathrm{T}] \mapsto [\mathrm{Hom}_{\mathrm{A}}(\Omega, \mathrm{T})]\) 与 \([\mathrm{M}] \mapsto [\Omega \otimes_{\mathrm{A}} \mathrm{M}]\) 在内射维数有限的有限生成 A-模的同构类所成的集合与投射维数有限的有限生成 A-模的同构类所成的集合之间定义了互逆的双射，*并且在相应的范畴之间定义了拟逆的范畴等价*。

¶ 7) 设 A 为局部 Macaulay 环，M 为有限生成的 Macaulay A-模。用 \(r(M)\)（相应地，\(t(M)\)）表示 \(\kappa_{A}\)-向量空间 \(\kappa_{A} \otimes_{A} M\)（相应地，\(\operatorname{Ext}_{A}^{p}(\kappa_{A}, M)\)，其中 \(p = \operatorname{prof}(M)\)）的维数，参见 § 1 习题 16。

\begin{enumerate}
\def\labelenumi{\alph{enumi})}
\item
  假设 A 具有对偶化模 \(\Omega\)；设 \(M' = \operatorname{Ext}_A^c(M, \Omega)\)，其中 \(c = \dim(A) - \dim_A(M)\)。证明等式 \(r(M') = t(M)\) 与 \(t(M') = r(M)\)（借助 § 3 命题 7 化为 \(c = 0\) 的情形，再对一个极大割序列取商化为 A 为 Artin 环的情形，并利用 § 8 命题 6）。
\item
  由此推出有 \(t(\mathrm{M}_{\mathfrak{p}}) \leqslant t(\mathrm{M})\) 对每个 \(\mathfrak{p} \in \operatorname{Spec}(\mathrm{A})\) 成立（借助 § 2 习题 2 与 § 1 习题 16 化为 A 完备的情形，并应用 a)）。
\item
  为使模 M 是对偶化的，必须且只需它是忠实的且 \(t(\mathrm{M}) = 1\)（化为 A 完备的情形；应用 a) 与定理 1 的推论）。
\end{enumerate}

\begin{enumerate}
\def\labelenumi{\arabic{enumi})}
\setcounter{enumi}{7}
\tightlist
\item
  设 A 为局部 Gorenstein 环，a 为高度为零的非零理想，b 为 a 的零化子。
\end{enumerate}

\begin{enumerate}
\def\labelenumi{\alph{enumi})}
\item
  证明理想 b 的高度为零，且 A-模 A/b 没有嵌入的伴随素理想（注意与 A/b 相伴的每个素理想都与 A 相伴）。
\item
  假设 A-模 A/α 没有嵌入的伴随素理想；证明 α 是 β 的零化子（化为 dim(A) = 0 的情形，并利用 Matlis 对偶）。
\item
  为使 A/a 是 Macaulay 环，必须且只需 A/b 亦然（利用定理 1 的推论与命题 3 的推论）。当这些条件满足时，A/a-模 b 是对偶化的。
\end{enumerate}

\begin{enumerate}
\def\labelenumi{\arabic{enumi})}
\setcounter{enumi}{8}
\tightlist
\item
  设 \(\Gamma\) 为加法单群 \(\mathbf{N}\) 的一个子单群，使得 \(\mathbf{N} - \Gamma\) 有限；用 \(c\) 表示使 \(\Gamma\) 包含所有整数 \(\geqslant c\) 的最小整数。设 \(k\) 为域，B 为环 \(k[[T]]\)，\(\mathrm{K} = k((\mathrm{T}))\) 为其分式域，A 为 B 的一个子 \(k\)-代数，由元素 \(\mathrm{T}^{\gamma}\)（\(\gamma \in \Gamma\)）生成。
\end{enumerate}

\begin{enumerate}
\def\labelenumi{\alph{enumi})}
\item
  证明 A 是维数 1 的局部整 Noether 环，其分式域为 K，整闭包为 B；理想 c = A : B 等于 BT \(^{c}\)。
\item
  为使 A 是 Gorenstein 环，必须且只需 \(c = 2\operatorname{Card}(\mathbf{N} - \Gamma)\)。
\item
  设 \(\Omega\) 为 K 中由元素 \(T^{-\alpha}\)（\(\alpha \notin \Gamma\)）生成的 A-子模。证明 \(K/\Omega\) 是 Matlis A-模（利用 § 8 命题 2），且 \(\Omega\) 是对偶化 A-模。
\end{enumerate}

\begin{enumerate}
\def\labelenumi{\arabic{enumi})}
\setcounter{enumi}{9}
\tightlist
\item
  设 A 为 Noether 环。设 I 为一个有界内射 A-复形，其同调为有限型的；对每个 A-复形 C，用 D(C) 表示复形 Homgr \(_{A}\) (C,I)。
\end{enumerate}

\begin{enumerate}
\def\labelenumi{\alph{enumi})}
\tightlist
\item
  证明下列条件等价：
\end{enumerate}

\begin{enumerate}
\def\labelenumi{(\roman{enumi})}
\item
  对每个有限生成的 A-模 M，复形的典范态射 \(\boldsymbol{\alpha}_{\mathrm{M}}:\mathrm{M}\to\mathbf{D}(\mathbf{D}(\mathrm{M}))\)（第 5 节）是拟同构；
\item
  对每个同调为有限型的有界 A-复形 C，态射 \(\alpha_{\mathrm{C}}: \mathrm{C} \to \mathbf{D}(\mathbf{D}(\mathrm{C}))\)（见该处）是拟同构；
\item
  典范态射 \(\alpha_{\mathrm{A}}:\mathrm{A}\to \mathrm{Homgr}_{\mathrm{A}}(\mathrm{I},\mathrm{I})\) 是拟同构。
\end{enumerate}

（仿照定理 1 的证明。）

若复形 I 有界、内射，H(I) 有限生成，且满足上述等价条件，我们就说复形 I 是对偶化的。

\begin{enumerate}
\def\labelenumi{\alph{enumi})}
\setcounter{enumi}{1}
\item
  若 A 具有对偶化模 \(\Omega\)，则 \(\Omega\) 的任一有限长度的内射分解 (I, δ) 都定义一个对偶化复形（见该处）。
\item
  设 B 为一个作为有限生成 A-模的 A-代数，I 为一个 A-模的对偶化复形；证明 B-模复形 \(\mathrm{Homgr}_{\mathrm{A}}(\mathrm{B},\mathrm{I})\) 是对偶化的。因此，有限维 Gorenstein 环的每个商环都具有对偶化复形。
\item
  设 I 为有界内射 A-复形。为使 I 是对偶化的，必须且只需对 A 的每个极大理想 m，\(A_{m}\)-模复形 \(I_{m}\) 都是对偶化的。
\item
  设 I 为对偶化 A-复形，P 为秩 1 的投射 A-模，n 为整数。复形 \(1 \otimes_{\mathrm{A}} \mathrm{P}(n)\) 是对偶化的。
\item
  设 I, J 为同调有限型的有界内射 A-复形，\(u: I \rightarrow J\) 为拟同构；为使 J 是对偶化的，必须且只需 I 是对偶化的。
\end{enumerate}

¶ 11) 设 A 为 Noether 环。

\begin{enumerate}
\def\labelenumi{\alph{enumi})}
\tightlist
\item
  假设 \(\Lambda\) 是局部的。设 P, Q 为有界平坦 A-复形。假设 \(\mathrm{H}(\mathrm{P}\otimes_{\mathrm{A}}\mathrm{Q})\) 在次数 \(\neq 0\) 处为零，在次数 0 处是秩 1 的自由模。证明存在一个整数 \(p\in \mathbf{Z}\) 以及拟同构 \(\mathrm{A}\to \mathrm{P}(p)\) 与 \(\mathrm{A}\to \mathrm{Q}(-p)\)。
\end{enumerate}

（利用 A, X 第~66 页的引理 1 与第~62 页的命题 1，化为 P 与 Q 在右边为零的情形；利用该处的命题 1 与 II, § 5, 第 4 节的定理 3，然后构造复形 P' 与 Q' 使 P = A ⊕ P' 且 Q = A ⊕ Q'，并证明有 H(P') = H(Q') = 0。）

\begin{enumerate}
\def\labelenumi{\alph{enumi})}
\setcounter{enumi}{1}
\item
  设 I, J 为对偶化 A-复形（习题 10）。若 A 是局部的，证明存在一个整数 n 以及一个从 J 到 I(n) 上的拟同构（设 P = Homgr \(_A\) (I, J)，Q = Homgr \(_A\) (J, I)；利用第 7 节的态射 \(\nu\) 以及相对于对偶化复形 1 的拟同构 \(\alpha_{J}\) : J → Homgr \(_A\) (Q, I)，构造一个同态 P ⊗ \(_A\) Q → A，并应用 a)）。
\item
  在一般情形，证明存在一个整数 n、一个秩 1 的投射 A-模 L，以及一个从 J 到 \(I \otimes_{A} L(n)\) 上的同态（设 \(L = H(\text{Homgr}_{A}(I, J))\)，并应用 b)）。
\end{enumerate}

\begin{enumerate}
\def\labelenumi{\arabic{enumi})}
\setcounter{enumi}{11}
\tightlist
\item
  设 \(k\) 为域，R 为一个 \(k\)-分次代数，型为 \(\mathbf{N}\)，满足 \(\mathrm{R}_0 = k\) 且 R 是维数 \(d\) 的 Macaulay 环。对每个有限生成的分次 R-模 M，用 \(\mathrm{P}_{\mathrm{M}}(\mathrm{T})\) 表示 Poincaré 级数 \(\sum_{i} (\dim_k \mathrm{M}_i) \mathrm{T}^i\)，并令 \(\mathrm{Q}_{\mathrm{M}}(\mathrm{T}) = (1 - \mathrm{T})^{\dim(M)} \mathrm{P}_{\mathrm{M}}(\mathrm{T})\)；它是 \(\mathbf{Z}[\mathrm{T}, \mathrm{T}^{-1}]\) 中的元素（VIII, § 6, 第 3 节, 命题 5）。
\end{enumerate}

\begin{enumerate}
\def\labelenumi{\alph{enumi})}
\item
  设 \(\Omega\) 为对偶化 R-模；证明 \(\Omega\) 具有一个分次 R-模结构，使得 \(Q_{\Omega}(T) = (-1)^{d} Q_{\mathbb{R}}(T^{-1})\)（把 R 写成多项式代数 \(A = k[X_1, \ldots, X_n]\) 的商，其中 \(\deg(X_i) = d_i\)；若 L 是 A-模 R 的一个分次自由分解，注意复形 \(\mathrm{Homgr}_{\mathrm{A}}(\mathrm{L}, \mathrm{A})(-d)\) 定义了 \(\Omega\) 的一个这样的结构，并利用 VIII, § 4, 第 2 节的例 3）。
\item
  假设 R 是 Gorenstein 环；证明存在整数 \(a \in \mathbf{Z}\) 满足 \(\mathrm{Q}_{\mathbb{R}}(\mathrm{T}^{-1}) = (-1)^{d}\mathrm{T}^{a}\mathrm{Q}_{\mathbb{R}}(\mathrm{T})\)。
\item
  假设 R 是整环，且存在整数 \(a \in \mathbf{Z}\) 使得 \(\mathrm{Q}_{\mathrm{R}}(\mathrm{T}^{-1}) = (-1)^{d}\mathrm{T}^{a}\mathrm{Q}_{\mathrm{R}}(\mathrm{T})\)。证明 R 是 Gorenstein 环（赋予 \(\Omega\) 一个使 \(\mathrm{Q}_{\Omega} = \mathrm{Q}_{\mathrm{R}}\) 的分次结构，并证明 \(\Omega_{0}\) 的一个非零元素构成 \(\Omega\) 的基）。
\item
  设 R 为分次 k-代数 \(k[X,Y]/(X^{3},XY,Y^{2})\)；证明有 \(Q_{R}(T^{-1}) = T^{-2}Q_{R}(T')\)，但 R 不是 Gorenstein 环。
\end{enumerate}

\exercisesection{局部上同调，Grothendieck 对偶}

\begin{enumerate}
\def\labelenumi{\arabic{enumi})}
\tightlist
\item
  设 A 为环，J 为 A 的理想，M 为 A-模。用 \(H_{A,J}(M)\) 或简作 \(H_{J}(M)\) 表示分次 A-模 \(\varinjlim_{n}\operatorname{Ext}_{A}(A/J^{n},M)\)。若 A 是局部 Noether 环，则分次 A-模 \(H_{A,m_{A}}(M)\) 与 \(H_{A}(M)\) 等同。
\end{enumerate}

\begin{enumerate}
\def\labelenumi{\alph{enumi})}
\item
  对 \(H_{J}(M)\) 建立与 \(H_{A}(M)\) 关于 A-模的同态与正合序列相类似的性质（第 1 节）。
\item
  对每个 \(i \geqslant 1\) 定义同构 \(\mathrm{H}_{\mathrm{J}}^{i+1}(\mathrm{M}) \to \varinjlim \operatorname{Ext}_{\mathrm{A}}^{i}(\mathrm{J}^{n}, \mathrm{M})\)。
\item
  假设环 A 是维数 d 的局部 Noether 环，且理想 J 由一个完全割序列 \((x_{1},\ldots,x_{r})\) 生成。构造一个同构

\[
\tau^ {i} (\mathrm{J}; \mathrm{M}): ^ {\prime} \operatorname{Tor} _ {r - i} ^ {\mathrm{A}} (\mathrm{M}, \mathrm{H} _ {\mathrm{J}} ^ {d} (\mathrm{A})) \longrightarrow \mathrm{H} _ {\mathrm{J}} ^ {i} (\mathrm{M})
\]

它推广了第 2 节的同构 \(\tau^{i}(\mathrm{M})\)；特别地，有 \(\mathrm{H}_{\mathrm{J}}^{i}(\mathrm{A}) = 0\) 对 \(i > r\) 成立。

\item
  对 A 的每个包含于 J 中的理想 K，用 \(\rho_{\mathrm{K},\mathrm{J}}:\mathrm{H}_{\mathrm{J}}(\mathrm{M})\to \mathrm{H}_{\mathrm{K}}(\mathrm{M})\) 表示由 \(\mathrm{A / K^n}\) 到 \(\mathrm{A / J^n}\) 的典范映射导出的同态。设 I 为 A 的一个理想；构造 A-模的一个长正合序列

\[
\dots \longrightarrow \mathrm{H} _ {\mathrm{I} +, \mathrm{J}} ^ {i} (\mathrm{M}) \stackrel {\alpha} {\longrightarrow} \mathrm{H} _ {\mathrm{I}} ^ {i} (\mathrm{M}) \oplus \mathrm{H} _ {\mathrm{J}} ^ {i} (\mathrm{M}) \stackrel {\beta} {\longrightarrow} \mathrm{H} _ {\mathrm{I} \cap \mathrm{J}} ^ {i} (\mathrm{M}) \longrightarrow \mathrm{H} _ {\mathrm{I} \mid \mathrm{J}} ^ {i + 1} (\mathrm{M}) \longrightarrow \dots ,
\]

\[
\text {   où   } \alpha (x) = (\rho_ {\mathrm{I,I+J}} (x), \rho_ {\mathrm{J,I+J}} (x)) \text {   et   } \beta (x, y) = \rho_ {\mathrm{I} \cap \mathrm{J}, \mathrm{I}} (x) - \rho_ {\mathrm{I} \cap \mathrm{J}, \mathrm{J}} (y).
\]

\end{enumerate}

\begin{enumerate}
\def\labelenumi{\arabic{enumi})}
\setcounter{enumi}{1}
\tightlist
\item
  设 \(\Lambda\) 为维数 \(d\) 的局部 Macaulay 环。证明为使 A 是 Gorenstein 环，必须且只需 \(\Lambda\)-模 \(\mathrm{H}_{\mathrm{A}}^{d}(\Lambda)\) 是内射的。
\end{enumerate}

¶ 3) 设 \(\rho: A \to B\) 为局部 Noether 环之间的局部同态，使得 \(\rho(\mathfrak{m}_{A})B\) 是 B 的定义理想。

\begin{enumerate}
\def\labelenumi{\alph{enumi})}
\item
  对每个 B-模 N，构造一个从 \(H_{A}(N)\) 到 \(H_{B}(N)\) 上的自然 A-线性同构。（利用 N 的一个内射分解，证明只需证明有 \(H_{A}^{i}(J)=0\) 对所有 i\textgreater0 与每个不可分解内射 B-模 J 成立。设 q 为 \(\operatorname{Ass}_{\mathrm{B}}(\mathrm{J})\) 中唯一的元素，并设 \(p=\rho^{-1}(q)\)。若 \(p\neq m_{A}\)，注意存在 \(a\in m_{A}\) 使得乘法映射 \(a_{J}\) 是双射。若 \(p=m_{A}\)，且设 I 为 A-模 J 的一个极小内射分解，注意对每个 n 有 \(\operatorname{Ass}_{\mathrm{A}}(\mathrm{I}^{n})=\{\mathfrak{m}_{\mathrm{A}}\}\)。）
\item
  进而假设 B 是平坦 A-模。对每个 A-模 M，构造一个从 \(\mathrm{B}\otimes_{\mathrm{A}}\mathrm{H}_{\mathrm{A}}(\mathrm{M})\) 到 \(\mathrm{H}_{\mathrm{B}}(\mathrm{B}\otimes_{\mathrm{A}}\mathrm{M})\) 上的自然 B-线性同构。
\end{enumerate}

\begin{enumerate}
\def\labelenumi{\arabic{enumi})}
\setcounter{enumi}{3}
\item
  设 A 为局部环，M 为非零的有限生成 A-模，维数为 e。证明 \(\mathrm{H}_{\mathrm{A}}^{e}(\mathrm{M})\) 不为零（化为环 A 完备的情形，再借助习题 3 化为 A 正则的情形）。
\item
  设 I 为 N 的有限子集。
\end{enumerate}

\begin{enumerate}
\def\labelenumi{\alph{enumi})}
\item
  构造一个正则局部环 B 与一个 B-模 N，使得满足 \(H_{\mathrm{B}}^{p}(\mathrm{N}) \neq 0\) 的整数 p 所成的集合等于 I（考虑直和）。
\item
  构造一个局部环 A，使得满足 \(\mathrm{H}_{\mathrm{A}}^{p}(\mathrm{A}) \neq 0\) 的整数 p 所成的集合等于 I（取 \(A = B \oplus N\)，其中 N 是平方为零的理想，并利用习题 3）。
\end{enumerate}

\begin{enumerate}
\def\labelenumi{\arabic{enumi})}
\setcounter{enumi}{5}
\tightlist
\item
  设 A 为具有对偶化模的局部 Noether 环，M 为有限生成 A-模，n 为整数。
\end{enumerate}

\begin{enumerate}
\def\labelenumi{\alph{enumi})}
\item
  为使 A-模 \(H_{\mathrm{A}}^{n}(\mathrm{M})\) 的长度有限，必须且只需对 A 的每个异于 \(m_{A}\) 的素理想 p，有 \(\mathrm{H}_{\mathrm{A}_{\mathfrak{p}}}^{n-\dim(\mathrm{A}/\mathfrak{p})}(\mathrm{M}_{\mathfrak{p}})=0\)（利用 Grothendieck 对偶）。
\item
  为使 \(\mathrm{H}_{\mathrm{A}}^{i}(\mathrm{M})\) 对每个 \(i \leqslant n\) 长度有限，必须且只需有 \(\operatorname{prof}(\mathrm{M}_{\mathfrak{p}}) + \dim(\mathrm{A}/\mathfrak{p}) > n\) 对每个异于 \(m_{A}\) 的素理想 p 成立
\end{enumerate}

\begin{enumerate}
\def\labelenumi{\arabic{enumi})}
\setcounter{enumi}{6}
\tightlist
\item
  设 A 为环，M 为 \(\Lambda\)-模，\(\mathbf{x} = (x_i)_{i \in I}\) 为 A 中元素的有限族。对每个整数 \(n \geqslant 1\)，用 \(\mathbf{x}^n\) 表示族 \((x_i^n)_{i \in I}\)。设 \(\Delta(\mathbf{x})\) 为 \(\mathrm{A}^{\mathrm{I}}\) 的一个自同态（由 \(\Delta(\mathbf{x})(e_i) = x_i e_i\) 定义）。
\end{enumerate}

\begin{enumerate}
\def\labelenumi{\alph{enumi})}
\tightlist
\item
  对每对整数 \((n, m)\)（满足 \(n \geqslant m\)），证明映射
\end{enumerate}

\[
\operatorname{Homgr} \left(\Lambda \left(\Delta \left(\mathbf {x} ^ {n - m}\right)\right), 1 _ {\mathrm{M}}\right): \mathbf {K} ^ {\bullet} \left(\mathbf {x} ^ {m}, \mathrm{M}\right) \longrightarrow \mathbf {K} ^ {\bullet} \left(\mathbf {x} ^ {n}, \mathrm{M}\right)
\]

定义复形的一个归纳系统；用 \(\mathbf{K}^{\bullet}(\mathbf{x}^{\infty},\mathrm{M})\) 表示这个系统的极限，用 \(\mathrm{H}^{\bullet}(\mathbf{x}^{\infty},\mathrm{M})\) 表示它的上同调。

\begin{enumerate}
\def\labelenumi{\alph{enumi})}
\setcounter{enumi}{1}
\tightlist
\item
  在 I 上选取一个全序。对 I 的每个子集 J，令 \(x_{J} = \prod_{j \in J} x_{j}\)。
\end{enumerate}

对 \(p \in \mathbf{Z}\)，令 \(\mathcal{C}^p = \bigoplus_{|J| = p} A_{x_J}\)；对 \(a \in A_{x_J}\)（其中 \(\text{Card}(J) = p\)），令 \(d^p(a) = \sum_{i \notin J} (-1)^{\varepsilon(J, i)} \frac{a}{1}\)，这里取值于 \(A_{x_{J \cup \{i\}}}\) 中，其中 \(\varepsilon(J, i)\) 是 \(J\) 中严格小于 \(i\) 的元素的个数。证明 \((\mathcal{C}, d)\) 是一个复形，并定义一个从 \(\mathbf{K}^\bullet(\mathbf{x}^\infty, M)\) 到复形 \(\mathcal{C} \otimes_A M\) 上的典范同构。

8) 沿用上一习题的记号，并进而假设环 A 是 Noether 的。用 x 表示由 x 生成的理想。

\begin{enumerate}
\def\labelenumi{\alph{enumi})}
\item
  设 m 为整数。证明存在整数 \(n_{0} \geqslant m\)，使得当 \(n \geqslant n_{0}\) 时，态射 \(\boldsymbol{\Lambda}(\Delta(\mathbf{x})^{n-m}) : \mathbf{K}_{\bullet}(\mathbf{x}^{n}, \mathrm{A}) \longrightarrow \mathbf{K}_{\bullet}(\mathbf{x}^{m}, \mathrm{A})\) 在次数 \(\geqslant 1\) 的同调上诱导 0（注意有 \(\boldsymbol{\Lambda}^{p}(\Delta(\mathbf{x}^{n-m}))(\mathbf{K}_{p}(\mathbf{x}^{n}, \mathrm{A})) \subset x^{n-m}\mathbf{K}_{p}(\mathbf{x}^{m}, \mathrm{A})\)；由 Artin–Rees 引理推出存在整数 \(n_{0}\)，使 \(Z_{p}(\mathbf{x}^{m}, \mathrm{A}) \cap x^{n_{0}}\mathbf{K}_{\bullet}(\mathbf{x}^{m}, \mathrm{A})\) 包含于 \(x^{rm}Z_{p}(\mathbf{x}^{m}, \mathrm{A})\)，从而包含于 \(B_{p}(\mathbf{x}^{m}, \mathrm{A})\)）。
\item
  定义一个从 \(\varinjlim_{n}\mathrm{Ext}_{\mathrm{A}}^{p}(\mathrm{A}/x^{n},\mathrm{M})\) 到 \(\mathrm{H}^p (\mathbf{x}^\infty ,\mathrm{M})\) 上的典范同构（化为证明 \(\mathrm{H}^p (\mathbf{x}^\infty ,\mathrm{M})\) 在 \(p\geqslant 1\) 且 M 内射时为零；在这一情形注意 \(\mathrm{H}^p (\mathbf{x}^\infty ,\mathrm{M})\) 与 \(\mathrm{Hom}_{\mathrm{A}}(\mathrm{H}_p(\mathbf{x}^n,\mathrm{M}))\) 的归纳极限等同，并利用 \(a)\)）。特别地，若 A 是局部的且 \(x\) 是 A 的定义理想，则分次 A-模 \(\mathrm{H}_{\mathrm{A}}(\mathrm{M})\) 典范同构于 \(\mathrm{H}^{\bullet}(\mathbf{x}^{\infty},\mathrm{M})\)。
\item
  假设环 A 是维数 \(d\) 的局部 Noether 环；设 \(\mathbf{x} = (x_{1},\ldots ,x_{d})\) 为 \(\mathfrak{m}_{\Lambda}\) 中元素的一个极大割序列。证明 A-模 \(\mathrm{H}_{\mathrm{A}}^{d}(\mathrm{M})\) 与 A-模的归纳系统 \((\mathrm{M}_p,u_{qp})\) 的极限等同，其中 \(\mathrm{M}_p\) 为商模 \(\mathrm{M} / (x_1^p\mathrm{M} + \dots +x_d^p\mathrm{M})\)，而 \(u_{qp}:\mathrm{M}_p\to \mathrm{M}_q\)（\(q\geqslant p\)）是由比率为 \((x_{1}\dots x_{d})^{q - p}\) 的乘法映射诱导的同态。
\end{enumerate}

\begin{enumerate}
\def\labelenumi{\arabic{enumi})}
\setcounter{enumi}{8}
\tightlist
\item
  设 A 为 Noether 局部环，d 为其维数，\(\mathbf{x} = (x_{1}, \ldots, x_{d})\) 为 \(m_{A}\) 中元素的一个极大割序列。
\end{enumerate}

\begin{enumerate}
\def\labelenumi{\alph{enumi})}
\item
  证明存在整数 \(m_0\)，使得当 \(m \geqslant m_0\) 且 \(n \geqslant 0\) 时，有 \((x_1^m \ldots x_d^m)^n \not\in (x_1^{m(n+1)}, \ldots, x_d^{m(n+1)})\)（利用习题 4 与 7, b)）。
\item
  假设 \(\Lambda\) 包含一个特征为 \(p\) 的域。证明有 \((x_{1} \ldots x_{d})^{n} \notin (x_{1}^{n+1}, \ldots, x_{d}^{n+1})\) 对每个 \(n \geqslant 0\) 成立，即 A 满足条件 (CM)（§ 2 习题 3；考察自同态 \(x \mapsto x^p\)）。因此，每个包含域的 Noether 局部环都满足条件 (CM) \(^{1}\)。
\item
  再假设 A 是正则的；设 \(\rho: A \to B\) 为一个单的环同态，使 B 成为有限生成的 A-模。证明 A-模 \(\rho(A)\) 是 B 的直因子（参见 § 2, 习题 3）；化为 B 是整环从而是局部的情形，并应用 § 4 习题 8）。
\end{enumerate}

\begin{enumerate}
\def\labelenumi{\arabic{enumi})}
\setcounter{enumi}{9}
\tightlist
\item
  设 \(k\) 为域，A 为 \(k\) 上的维数 \(d\) 的正则局部代数，且扩张 \(k \to \kappa_{\mathrm{A}}\) 是平凡的。A-模 \(\mathrm{H}_{\mathrm{A}}^{d}(\mathrm{A})\) 与 \(\mathrm{A}' = \mathrm{Hom}_k^{cont}(\mathrm{A}, k)\) 都是 Matlis 模，从而同构；我们要在这两个模之间把一个同构明显地写出来。
\end{enumerate}

\begin{enumerate}
\def\labelenumi{\alph{enumi})}
\item
  设 \(\mathbf{x} = (x_{1}, \ldots, x_{d})\) 为 A 的一个坐标系。对每个整数 \(p \geqslant 0\)，用 \(A_{p}\) 表示 k-代数 \(\mathrm{A}/(x_{1}^{p}, \ldots, x_{d}^{p})\)。单项式 \(x_{1}^{a_{1}} \ldots x_{d}^{a_{d}}\)（其中 \(0 \leqslant a_{i} < p\)）的类构成 \(A_{p}\) 在 k 上的一组基；用 \(\rho_{p}\) 表示 \(A_{p}\) 上的线性形式，它在 \(x_{1}^{p-1} \ldots x_{d}^{p-1}\) 的类上取值 1，在基的其余元素上取值 0。证明 k-线性映射 \(\tilde{\rho}_{p}: A_{p} \to \operatorname{Hom}_{k}(A_{p}, k)\)（它由双线性形式 \((a, b) \mapsto \rho_{p}(ab)\) 导出）是同构（可直接论证，或利用 § 5 习题 10）。
\item
  对 \(q \geqslant p\)，设 \(u_{qp}: A_p \to A_q\) 为由比率为 \((x_1 \ldots x_d)^{q-p}\) 的乘法映射诱导的 A-线性同态，并设 \(\pi_{pq}: A_q \to A_p\) 为商映射；证明有 \(\tilde{\rho}_q \circ u_{qp} = {}^t \pi_{pq} \circ \tilde{\rho}_p\)。取归纳极限并利用习题 8, c)，由此得到一个 A-线性同构 \(\tilde{\rho}: H_A^d(A) \to A'\)。
\item
  假设 \(d=1\)，于是 A 是离散赋值环，且 \(x_{1}=x\) 是 A 的一个单值化元。证明单值化元 \(x\) 的选取使 \(\mathrm{H}_{\mathrm{A}}^{1}(\mathrm{A})\) 得以与 K/A 等同；同构 \(\rho:\mathrm{K}/\mathrm{A}\to\mathrm{A}'\) 于是由 \(\rho(\varphi)(a)=\mathrm{Res}(a\varphi)\) 定义，其中 \(\mathrm{Res}:\mathrm{K}/\mathrm{A}\to k\) 是到因子 \(k x^{-1}\)（在典范分解 \(\mathrm{K}/\Lambda=\bigoplus_{n\geqslant 1}k x^{-n}\) 中）上的投射。
\end{enumerate}

\begin{enumerate}
\def\labelenumi{\arabic{enumi})}
\setcounter{enumi}{10}
\item
  设 A 为 Noether 局部环，I 为 A 上的一个 Matlis 模，M 为有限生成 A-模。令 \(\mu^{p} = \dim_{\kappa_{A}} \operatorname{Ext}_{A}^{p}(\kappa_{A}, M)\) 对每个 \(p \geqslant 0\)（参见 § 8, 习题 4）。证明分次 A-模 H \(_{A}\) (M) 与一个复形 \(\ldots \to 0 \to I^{\mu^{0}} \to I^{\mu^{1}} \to \ldots\) 的同调等同（见该处）。特别地，若 M 非零，则有 \(\mu^{p} \neq 0\) 对 \(p = \text{prof}(M)\) 与 \(p = \dim(M)\)（习题 4）。
\item
  设 A 为 Noether 局部环，M 为维数 d 的有限生成 A-模；假设 M 是 Buchsbaum 模（§ 2, 习题 7）。
\end{enumerate}

\begin{enumerate}
\def\labelenumi{\alph{enumi})}
\item
  证明 \(\mathrm{H}_{\mathrm{A}}^{i}(\mathrm{M})\) 被 \(m_{A}\) 零化（对 \(i \neq d\)）（先处理 d = 1 的情形，再对 d 作归纳）。
\item
  设 \(x\) 为关于 \(\mathfrak{m}\) 的一个割元素；对 \(0 \leqslant i \leqslant d - 2\) 定义正合序列 \(0 \to \mathrm{H}_{\mathrm{A}}^{i}(\mathrm{M}) \to \mathrm{H}_{\mathrm{A}}^{i}(\mathrm{M}/x\mathrm{M}) \to \mathrm{H}_{\mathrm{A}}^{i+1}(\mathrm{M}) \to 0\)。
\item
  证明公式 \(i(\mathrm{M}) = \sum_{i=0}^{d-1} \left( \begin{array}{cc} d & 1 \\ i & \end{array} \right) \lg(\mathrm{H}_{\mathrm{A}}^i(\mathrm{M}))\)（对整数 \(d \geqslant 1\) 作归纳，利用 \(b\) ) 与 § 2 习题 8, \(d\)）。
\end{enumerate}

\specialchapter{记号索引}

\(\mathrm{prof}_{\mathrm{A}}(\mathrm{J};\mathrm{M})\)、\(\mathrm{prof}(\mathrm{J};\mathrm{M})\)、\(\mathrm{prof}_{\mathrm{A}}(\mathrm{M})\)、\(\mathrm{prof}(\mathrm{M})\)：第~2 页。\(\mathbf{K}^{\bullet}(\mathbf{x},\mathbf{M})\)：第~5 页。\(\mathrm{H}^{\bullet}(\mathbf{x},\mathbf{M})\) 第~5 页。\(\mathrm{prof}_{\mathrm{F}}(\mathrm{M})\) 第~10 页。\(\mathrm{gradc}(\mathrm{N})\) 第~11 页。\(\mathrm{dp}_{\mathrm{A}}(\mathrm{M})\)、\(\mathrm{di}_{\mathrm{A}}(\mathrm{M})\) 第~37 页。\(\mathrm{dh(A)}\) 第~38 页。\(\delta (\Lambda)\) 第~60 页。\(\alpha_{\mathrm{J}},\beta_{\mathrm{J}}\) 第~63 页。\(\Omega_k(\mathrm{A}),d_{\mathrm{A / k}}\) 第~78 页。\(\Omega (\mathrm{A})\) 第~93 页。\(\mathcal{I}(\mathrm{A})\) 第~105 页。\\
I(p), \(e_{\mathfrak{p}}\) 第~105 页。\(\mathrm{D_A(M),D_A(f),}\alpha_{\mathrm{M}}\) 第~112 页。\(\theta (\mathrm{M},\mathrm{P}),\rho (\mathrm{M},\mathrm{P})\) 第~121 页。\(\mathbf{D}(\mathbf{C})\)，第~134 页。\(\mathbf{D}_{\mathrm{A}}(f),\mathbf{D}(\mathbf{M})\) 第~135 页。\(\alpha_{\mathrm{M}}\) 第~136 页。\(\mathrm{H}_{\mathrm{A}}(\mathrm{M}),\mathrm{H}_{\mathrm{A}}(f)\) 第~142 页。\(\mathcal{D}\) 第~142 页。\(\mathcal{D}_{cs}\) 第~145 页。\(\tau^i (\mathbf{M})\) 第~146 页。\(\gamma^i (\mathbf{M}),\delta^i (\mathbf{M})\) 第~148 页。\(S_{k}\) 第~151 页, 习题 8。\(t_\mathrm{A}(\mathbf{M}),t(\mathbf{M})\) 第~153 页, 习题 16。\(\chi (\mathbf{x},\mathbf{N})\) 第~155 页, 习题 6。\(\chi (\mathbf{M})\) 第~157 页, 习题 2。\(rg(u),\mathfrak{d}(u)\) 第~159 页, 习题 8。\(h_i(\mathbf{A})\) 第~166 页, 习题 4。\(\mu^i (\mathfrak{p},\mathbf{M})\) 第~174 页, 习题 4。\(\mathrm{H}_{\mathrm{J}}(\mathrm{M})\) 第~179 页, 习题 1。\(\mathbf{K}^{\bullet}(\mathbf{x}^{\infty},\mathbf{M}),\mathrm{H}^{\bullet}(\mathbf{x}^{\infty},\mathbf{M})\) 第~180 页, 习题 7。

\specialchapter{术语索引}

绝对正则、正规（代数）：第~75 页。绝对正则、绝对正规代数：第~75 页。本质有限型代数：第~71 页。形式平展、形式非分歧代数：第~170 页, 习题 1 形式光滑代数：第~85 页。光滑代数：第~104 页。完全交环：第~65 页。Gorenstein 环：第~48 页。Grothendieck 环：第~169 页, 习题 5。局部整的环：第~16 页。Macaulay 环（或 Macaulay 环）：第~30 页。正规环：第~16 页。可表现环：第~58 页。正则环：第~55 页。Auslander-Buchsbaum（定理）：第~45 页。

Bass（定理）：第~49 页。Buchsbaum（模）：第~155 页, 习题 7。

Cohen（定理）：第~88 页。局部上同调（模）：第~142 页。完全割（理想）：第~62 页。对偶化复形：第~177 页, 习题 10 环的典范分支：第~15 页。

微差（理想）：第~167 页, 习题 7。环的同调维数：第~38 页。模的投射或内射维数：第~37 页。对偶化（复形）：第~177 页, 习题 10。对偶化（模）：第~125 页。Grothendieck 对偶：第~149 页。

本质有限型（代数）：第~71 页。本质生成（族）：第~71 页。

本质生成族：第~71 页。形式平展、形式非分歧（代数）：第~170 页, 习题 1。形式光滑（代数）：第~85 页。强割（子集）：第~28 页。

Gorenstein（环）：第~48 页。模的级：第~11 页。Grothendieck（环）：第~169 页, 习题 5。Grothendieck（对偶）：第~149 页。

Hartshorne（定理）：第~17 页。Hilbert-Burch（定理）：第~160 页, 习题 11。同调（维数）：第~38 页。完全割理想：第~62 页。内射（维数）：第~37 页。完全交（环）：第~65 页。光滑（形式光滑代数）：第~85 页。光滑（代数）：第~104 页。Macaulay（环）：第~30 页。Macaulay-Cohen（判别法）：第~31 页。Macaulay-Cohen（定理）：第~30 页。Macaulay（模）：第~30 页。Matlis（模）：第~110 页。Buchsbaum 模：第~155 页, 习题 7。局部上同调模：第~142 页。Matlis 模：第~110 页。对偶化模：第~125 页。Macaulay（或 Macaulay 的）模：第~23 页。完美模：第~158 页, 习题 4。纯模：第~27 页。正规（环）：第~16 页。完美（模）：第~158 页, 习题 4。强割子集：第~28 页。可表现（环）：第~58 页。模沿闭子集的深度：第~10 页。模、局部环的深度：第~2 页。投射（维数）：第~37 页。性质 \(S_k\)：第~151 页, 习题 8。素理想的 n 次符号幂：第~109 页。纯（模）：第~27 页。正则（环）：第~55 页。代数同态的提升：第~83 页。割（完全割理想）：第~62 页。Serre（定理）：第~20 页与第~54 页。模的底座：第~111 页。M-正则序列：第~8 页。Auslander-Buchsbaum 定理：第~45 页。Bass 定理：第~49 页。Cohen 定理：第~88 页。Hartshorne 定理：第~17 页。Hilbert-Burch 定理：第~160 页, 习题 11。Macaulay-Cohen 定理：第~30 页。Serre 定理：第~20 页与第~54 页。Zariski 定理：第~101 页与第~102 页。Zariski（定理）：第~101 页与第~102 页。
